% ####################################################################
% Basic configuration of Beamera class and Jagiellonian theme
% ####################################################################
\RequirePackage[l2tabu, orthodox]{nag}



\ifx\PresentationStyle\notset
  \def\PresentationStyle{dark}
\fi



% Options: t -- align text to the top of the frame.
\documentclass[10pt,t]{beamer}
\mode<presentation>
\usetheme[style=\PresentationStyle]{jagiellonian}





% ------------------------------------------------------------------------------------
% Procesing configuration files of Jagiellonian theme located in
% the directory "preambule".
% ------------------------------------------------------------------------------------
% Configuration for polish language
% Need description
\usepackage[polish]{babel}
% Need description
\usepackage[MeX]{polski}



% ------------------------------
% Better support of polish chars in technical parts of PDF
% ------------------------------
\hypersetup{pdfencoding=auto,psdextra}

% Package "textpos" give as enviroment "textblock" which is very usefull in
% arranging text on slides.

% This is standard configuration of "textpos"
\usepackage[overlay,absolute]{textpos}

% If you need to see bounds of "textblock's" comment line above and uncomment
% one below.

% Caution! When showboxes option is on significant ammunt of space is add
% to the top of textblock and as such, everyting put in them gone down.
% We need to check how to remove this bug.

% \usepackage[showboxes,overlay,absolute]{textpos}



% Setting scale length for package "textpos"
\setlength{\TPHorizModule}{10mm}
\setlength{\TPVertModule}{\TPHorizModule}


% ---------------------------------------
% Packages written for lectures "Geometria 3D dla twórców gier wideo"
% ---------------------------------------
% \usepackage{./Geometry3DPackages/Geometry3D}
% \usepackage{./Geometry3DPackages/Geometry3DIndices}
% \usepackage{./Geometry3DPackages/Geometry3DTikZStyle}
% \usepackage{./ProgramowanieSymulacjiFizykiPaczki/ProgramowanieSymulacjiFizykiTikZStyle}
% \usepackage{./Geometry3DPackages/mathcommands}


% ---------------------------------------
% Importing TikZ's libraries
% ---------------------------------------
\usetikzlibrary{arrows.meta}
\usetikzlibrary{positioning}





% % Configuration package "bm" that need for making bold symbols
\newcommand{\bmmax}{0}
\newcommand{\hmmax}{0}
\usepackage{bm}




% ---------------------------------------
% Packages for scientific texts
% ---------------------------------------
% \let\lll\undefined  % Sometimes you must use this line to allow
% "amsmath" package to works with packages with packages for polish
% languge imported
% /preambul/LanguageSettings/JagiellonianPolishLanguageSettings.tex.
% This comments (probably) removes polish letter Ł.
\usepackage{amsmath}  % Packages from American Mathematical Society (AMS)
\usepackage{amssymb}
\usepackage{amscd}
\usepackage{amsthm}
\usepackage{siunitx}  % Package for typsetting SI units.
\usepackage{upgreek}  % Better looking greek letters.
% Example of using upgreek: pi = \uppi


\usepackage{calrsfs}  % Zmienia czcionkę kaligraficzną w \mathcal
% na ładniejszą. Może w innych miejscach robi to samo, ale o tym nic
% nie wiem.







% !!!!!!!!!!!!!!!!!!!!!!!!!!!!!!
% !!!!!!!!!!!!!!!!!!!!!!!!!!!!!!
% EVIL STUFF
\if\JUlogotitle1
\edef\LogoJUPath{LogoJU_\JUlogoLang/LogoJU_\JUlogoShape_\JUlogoColor.pdf}
\titlegraphic{\hfill\includegraphics[scale=0.22]
{./JagiellonianPictures/\LogoJUPath}}
\fi

% ---------------------------------------
% Commands for handling colors
% ---------------------------------------


% Command for setting normal text color for some text in math modestyle
% Text color depend on used style of Jagiellonian

% Beamer version of command
\newcommand{\TextWithNormalTextColor}[1]{%
  {\color{jNormalTextFGColor}
    \setbeamercolor{math text}{fg=jNormalTextFGColor} {#1}}
}

% Article and similar classes version of command
% \newcommand{\TextWithNormalTextColor}[1]{%
%   {\color{jNormalTextsFGColor} {#1}}
% }



% Beamer version of command
\newcommand{\NormalTextInMathMode}[1]{%
  {\color{jNormalTextFGColor}
    \setbeamercolor{math text}{fg=jNormalTextFGColor} \text{#1}}
}


% Article and similar classes version of command
% \newcommand{\NormalTextInMathMode}[1]{%
%   {\color{jNormalTextsFGColor} \text{#1}}
% }




% Command that sets color of some mathematical text to the same color
% that has normal text in header (?)

% Beamer version of the command
\newcommand{\MathTextFrametitleFGColor}[1]{%
  {\color{jFrametitleFGColor}
    \setbeamercolor{math text}{fg=jFrametitleFGColor} #1}
}

% Article and similar classes version of the command
% \newcommand{\MathTextWhiteColor}[1]{{\color{jFrametitleFGColor} #1}}





% Command for setting color of alert text for some text in math modestyle

% Beamer version of the command
\newcommand{\MathTextAlertColor}[1]{%
  {\color{jOrange} \setbeamercolor{math text}{fg=jOrange} #1}
}

% Article and similar classes version of the command
% \newcommand{\MathTextAlertColor}[1]{{\color{jOrange} #1}}





% Command that allow you to sets chosen color as the color of some text into
% math mode. Due to some nuances in the way that Beamer handle colors
% it not work in all cases. We hope that in the future we will improve it.

% Beamer version of the command
\newcommand{\SetMathTextColor}[2]{%
  {\color{#1} \setbeamercolor{math text}{fg=#1} #2}
}


% Article and similar classes version of the command
% \newcommand{\SetMathTextColor}[2]{{\color{#1} #2}}










% ---------------------------------------
% Commands for few special slides
% ---------------------------------------
\newcommand{\EndingSlide}[1]{%
  \begin{frame}[standout]

    \begingroup

    \color{jFrametitleFGColor}

    #1

    \endgroup

  \end{frame}
}










% ---------------------------------------
% Commands for setting background pictures for some slides
% ---------------------------------------
\newcommand{\TitleBackgroundPicture}
{./JagiellonianPictures/Backgrounds/LajkonikDark.png}
\newcommand{\SectionBackgroundPicture}
{./JagiellonianPictures/Backgrounds/LajkonikLight.png}



\newcommand{\TitleSlideWithPicture}{%
  \begingroup

  \usebackgroundtemplate{%
    \includegraphics[height=\paperheight]{\TitleBackgroundPicture}}

  \maketitle

  \endgroup
}





\newcommand{\SectionSlideWithPicture}[1]{%
  \begingroup

  \usebackgroundtemplate{%
    \includegraphics[height=\paperheight]{\SectionBackgroundPicture}}

  \setbeamercolor{titlelike}{fg=normal text.fg}

  \section{#1}

  \endgroup
}










% ---------------------------------------
% Commands for lectures "Geometria 3D dla twórców gier wideo"
% Polish version
% ---------------------------------------
% Komendy teraz wykomentowane były potrzebne, gdy loga były na niebieskim
% tle, nie na białym. A są na białym bo tego chcieli w biurze projektu.
% \newcommand{\FundingLogoWhitePicturePL}
% {./PresentationPictures/CommonPictures/logotypFundusze_biale_bez_tla2.pdf}
\newcommand{\FundingLogoColorPicturePL}
{./PresentationPictures/CommonPictures/European_Funds_color_PL.pdf}
% \newcommand{\EULogoWhitePicturePL}
% {./PresentationPictures/CommonPictures/logotypUE_biale_bez_tla2.pdf}
\newcommand{\EUSocialFundLogoColorPicturePL}
{./PresentationPictures/CommonPictures/EU_Social_Fund_color_PL.pdf}
% \newcommand{\ZintegrUJLogoWhitePicturePL}
% {./PresentationPictures/CommonPictures/zintegruj-logo-white.pdf}
\newcommand{\ZintegrUJLogoColorPicturePL}
{./PresentationPictures/CommonPictures/ZintegrUJ_color.pdf}
\newcommand{\JULogoColorPicturePL}
{./JagiellonianPictures/LogoJU_PL/LogoJU_A_color.pdf}





\newcommand{\GeometryThreeDSpecialBeginningSlidePL}{%
  \begin{frame}[standout]

    \begin{textblock}{11}(1,0.7)

      \begin{flushleft}

        \mdseries

        \footnotesize

        \color{jFrametitleFGColor}

        Materiał powstał w ramach projektu współfinansowanego ze środków
        Unii Europejskiej w ramach Europejskiego Funduszu Społecznego
        POWR.03.05.00-00-Z309/17-00.

      \end{flushleft}

    \end{textblock}





    \begin{textblock}{10}(0,2.2)

      \tikz \fill[color=jBackgroundStyleLight] (0,0) rectangle (12.8,-1.5);

    \end{textblock}


    \begin{textblock}{3.2}(1,2.45)

      \includegraphics[scale=0.3]{\FundingLogoColorPicturePL}

    \end{textblock}


    \begin{textblock}{2.5}(3.7,2.5)

      \includegraphics[scale=0.2]{\JULogoColorPicturePL}

    \end{textblock}


    \begin{textblock}{2.5}(6,2.4)

      \includegraphics[scale=0.1]{\ZintegrUJLogoColorPicturePL}

    \end{textblock}


    \begin{textblock}{4.2}(8.4,2.6)

      \includegraphics[scale=0.3]{\EUSocialFundLogoColorPicturePL}

    \end{textblock}

  \end{frame}
}



\newcommand{\GeometryThreeDTwoSpecialBeginningSlidesPL}{%
  \begin{frame}[standout]

    \begin{textblock}{11}(1,0.7)

      \begin{flushleft}

        \mdseries

        \footnotesize

        \color{jFrametitleFGColor}

        Materiał powstał w ramach projektu współfinansowanego ze środków
        Unii Europejskiej w ramach Europejskiego Funduszu Społecznego
        POWR.03.05.00-00-Z309/17-00.

      \end{flushleft}

    \end{textblock}





    \begin{textblock}{10}(0,2.2)

      \tikz \fill[color=jBackgroundStyleLight] (0,0) rectangle (12.8,-1.5);

    \end{textblock}


    \begin{textblock}{3.2}(1,2.45)

      \includegraphics[scale=0.3]{\FundingLogoColorPicturePL}

    \end{textblock}


    \begin{textblock}{2.5}(3.7,2.5)

      \includegraphics[scale=0.2]{\JULogoColorPicturePL}

    \end{textblock}


    \begin{textblock}{2.5}(6,2.4)

      \includegraphics[scale=0.1]{\ZintegrUJLogoColorPicturePL}

    \end{textblock}


    \begin{textblock}{4.2}(8.4,2.6)

      \includegraphics[scale=0.3]{\EUSocialFundLogoColorPicturePL}

    \end{textblock}

  \end{frame}





  \TitleSlideWithPicture
}



\newcommand{\GeometryThreeDSpecialEndingSlidePL}{%
  \begin{frame}[standout]

    \begin{textblock}{11}(1,0.7)

      \begin{flushleft}

        \mdseries

        \footnotesize

        \color{jFrametitleFGColor}

        Materiał powstał w ramach projektu współfinansowanego ze środków
        Unii Europejskiej w~ramach Europejskiego Funduszu Społecznego
        POWR.03.05.00-00-Z309/17-00.

      \end{flushleft}

    \end{textblock}





    \begin{textblock}{10}(0,2.2)

      \tikz \fill[color=jBackgroundStyleLight] (0,0) rectangle (12.8,-1.5);

    \end{textblock}


    \begin{textblock}{3.2}(1,2.45)

      \includegraphics[scale=0.3]{\FundingLogoColorPicturePL}

    \end{textblock}


    \begin{textblock}{2.5}(3.7,2.5)

      \includegraphics[scale=0.2]{\JULogoColorPicturePL}

    \end{textblock}


    \begin{textblock}{2.5}(6,2.4)

      \includegraphics[scale=0.1]{\ZintegrUJLogoColorPicturePL}

    \end{textblock}


    \begin{textblock}{4.2}(8.4,2.6)

      \includegraphics[scale=0.3]{\EUSocialFundLogoColorPicturePL}

    \end{textblock}





    \begin{textblock}{11}(1,4)

      \begin{flushleft}

        \mdseries

        \footnotesize

        \RaggedRight

        \color{jFrametitleFGColor}

        Treść niniejszego wykładu jest udostępniona na~licencji
        Creative Commons (\textsc{cc}), z~uzna\-niem autorstwa
        (\textsc{by}) oraz udostępnianiem na tych samych warunkach
        (\textsc{sa}). Rysunki i~wy\-kresy zawarte w~wykładzie są
        autorstwa dr.~hab.~Pawła Węgrzyna et~al. i~są dostępne
        na tej samej licencji, o~ile nie wskazano inaczej.
        W~prezentacji wykorzystano temat Beamera Jagiellonian,
        oparty na~temacie Metropolis Matthiasa Vogelgesanga,
        dostępnym na licencji \LaTeX{} Project Public License~1.3c
        pod adresem: \colorhref{https://github.com/matze/mtheme}
        {https://github.com/matze/mtheme}.

        Projekt typograficzny: Iwona Grabska-Gradzińska \\
        Skład: Kamil Ziemian;
        Korekta: Wojciech Palacz \\
        Modele: Dariusz Frymus, Kamil Nowakowski \\
        Rysunki i~wykresy: Kamil Ziemian, Paweł Węgrzyn, Wojciech Palacz

      \end{flushleft}

    \end{textblock}

  \end{frame}
}



\newcommand{\GeometryThreeDTwoSpecialEndingSlidesPL}[1]{%
  \begin{frame}[standout]


    \begin{textblock}{11}(1,0.7)

      \begin{flushleft}

        \mdseries

        \footnotesize

        \color{jFrametitleFGColor}

        Materiał powstał w ramach projektu współfinansowanego ze środków
        Unii Europejskiej w~ramach Europejskiego Funduszu Społecznego
        POWR.03.05.00-00-Z309/17-00.

      \end{flushleft}

    \end{textblock}





    \begin{textblock}{10}(0,2.2)

      \tikz \fill[color=jBackgroundStyleLight] (0,0) rectangle (12.8,-1.5);

    \end{textblock}


    \begin{textblock}{3.2}(1,2.45)

      \includegraphics[scale=0.3]{\FundingLogoColorPicturePL}

    \end{textblock}


    \begin{textblock}{2.5}(3.7,2.5)

      \includegraphics[scale=0.2]{\JULogoColorPicturePL}

    \end{textblock}


    \begin{textblock}{2.5}(6,2.4)

      \includegraphics[scale=0.1]{\ZintegrUJLogoColorPicturePL}

    \end{textblock}


    \begin{textblock}{4.2}(8.4,2.6)

      \includegraphics[scale=0.3]{\EUSocialFundLogoColorPicturePL}

    \end{textblock}





    \begin{textblock}{11}(1,4)

      \begin{flushleft}

        \mdseries

        \footnotesize

        \RaggedRight

        \color{jFrametitleFGColor}

        Treść niniejszego wykładu jest udostępniona na~licencji
        Creative Commons (\textsc{cc}), z~uzna\-niem autorstwa
        (\textsc{by}) oraz udostępnianiem na tych samych warunkach
        (\textsc{sa}). Rysunki i~wy\-kresy zawarte w~wykładzie są
        autorstwa dr.~hab.~Pawła Węgrzyna et~al. i~są dostępne
        na tej samej licencji, o~ile nie wskazano inaczej.
        W~prezentacji wykorzystano temat Beamera Jagiellonian,
        oparty na~temacie Metropolis Matthiasa Vogelgesanga,
        dostępnym na licencji \LaTeX{} Project Public License~1.3c
        pod adresem: \colorhref{https://github.com/matze/mtheme}
        {https://github.com/matze/mtheme}.

        Projekt typograficzny: Iwona Grabska-Gradzińska \\
        Skład: Kamil Ziemian;
        Korekta: Wojciech Palacz \\
        Modele: Dariusz Frymus, Kamil Nowakowski \\
        Rysunki i~wykresy: Kamil Ziemian, Paweł Węgrzyn, Wojciech Palacz

      \end{flushleft}

    \end{textblock}

  \end{frame}





  \begin{frame}[standout]

    \begingroup

    \color{jFrametitleFGColor}

    #1

    \endgroup

  \end{frame}
}



\newcommand{\GeometryThreeDSpecialEndingSlideVideoPL}{%
  \begin{frame}[standout]

    \begin{textblock}{11}(1,0.7)

      \begin{flushleft}

        \mdseries

        \footnotesize

        \color{jFrametitleFGColor}

        Materiał powstał w ramach projektu współfinansowanego ze środków
        Unii Europejskiej w~ramach Europejskiego Funduszu Społecznego
        POWR.03.05.00-00-Z309/17-00.

      \end{flushleft}

    \end{textblock}





    \begin{textblock}{10}(0,2.2)

      \tikz \fill[color=jBackgroundStyleLight] (0,0) rectangle (12.8,-1.5);

    \end{textblock}


    \begin{textblock}{3.2}(1,2.45)

      \includegraphics[scale=0.3]{\FundingLogoColorPicturePL}

    \end{textblock}


    \begin{textblock}{2.5}(3.7,2.5)

      \includegraphics[scale=0.2]{\JULogoColorPicturePL}

    \end{textblock}


    \begin{textblock}{2.5}(6,2.4)

      \includegraphics[scale=0.1]{\ZintegrUJLogoColorPicturePL}

    \end{textblock}


    \begin{textblock}{4.2}(8.4,2.6)

      \includegraphics[scale=0.3]{\EUSocialFundLogoColorPicturePL}

    \end{textblock}





    \begin{textblock}{11}(1,4)

      \begin{flushleft}

        \mdseries

        \footnotesize

        \RaggedRight

        \color{jFrametitleFGColor}

        Treść niniejszego wykładu jest udostępniona na~licencji
        Creative Commons (\textsc{cc}), z~uzna\-niem autorstwa
        (\textsc{by}) oraz udostępnianiem na tych samych warunkach
        (\textsc{sa}). Rysunki i~wy\-kresy zawarte w~wykładzie są
        autorstwa dr.~hab.~Pawła Węgrzyna et~al. i~są dostępne
        na tej samej licencji, o~ile nie wskazano inaczej.
        W~prezentacji wykorzystano temat Beamera Jagiellonian,
        oparty na~temacie Metropolis Matthiasa Vogelgesanga,
        dostępnym na licencji \LaTeX{} Project Public License~1.3c
        pod adresem: \colorhref{https://github.com/matze/mtheme}
        {https://github.com/matze/mtheme}.

        Projekt typograficzny: Iwona Grabska-Gradzińska;
        Skład: Kamil Ziemian \\
        Korekta: Wojciech Palacz;
        Modele: Dariusz Frymus, Kamil Nowakowski \\
        Rysunki i~wykresy: Kamil Ziemian, Paweł Węgrzyn, Wojciech Palacz \\
        Montaż: Agencja Filmowa Film \& Television Production~-- Zbigniew
        Masklak

      \end{flushleft}

    \end{textblock}

  \end{frame}
}





\newcommand{\GeometryThreeDTwoSpecialEndingSlidesVideoPL}[1]{%
  \begin{frame}[standout]

    \begin{textblock}{11}(1,0.7)

      \begin{flushleft}

        \mdseries

        \footnotesize

        \color{jFrametitleFGColor}

        Materiał powstał w ramach projektu współfinansowanego ze środków
        Unii Europejskiej w~ramach Europejskiego Funduszu Społecznego
        POWR.03.05.00-00-Z309/17-00.

      \end{flushleft}

    \end{textblock}





    \begin{textblock}{10}(0,2.2)

      \tikz \fill[color=jBackgroundStyleLight] (0,0) rectangle (12.8,-1.5);

    \end{textblock}


    \begin{textblock}{3.2}(1,2.45)

      \includegraphics[scale=0.3]{\FundingLogoColorPicturePL}

    \end{textblock}


    \begin{textblock}{2.5}(3.7,2.5)

      \includegraphics[scale=0.2]{\JULogoColorPicturePL}

    \end{textblock}


    \begin{textblock}{2.5}(6,2.4)

      \includegraphics[scale=0.1]{\ZintegrUJLogoColorPicturePL}

    \end{textblock}


    \begin{textblock}{4.2}(8.4,2.6)

      \includegraphics[scale=0.3]{\EUSocialFundLogoColorPicturePL}

    \end{textblock}





    \begin{textblock}{11}(1,4)

      \begin{flushleft}

        \mdseries

        \footnotesize

        \RaggedRight

        \color{jFrametitleFGColor}

        Treść niniejszego wykładu jest udostępniona na~licencji
        Creative Commons (\textsc{cc}), z~uzna\-niem autorstwa
        (\textsc{by}) oraz udostępnianiem na tych samych warunkach
        (\textsc{sa}). Rysunki i~wy\-kresy zawarte w~wykładzie są
        autorstwa dr.~hab.~Pawła Węgrzyna et~al. i~są dostępne
        na tej samej licencji, o~ile nie wskazano inaczej.
        W~prezentacji wykorzystano temat Beamera Jagiellonian,
        oparty na~temacie Metropolis Matthiasa Vogelgesanga,
        dostępnym na licencji \LaTeX{} Project Public License~1.3c
        pod adresem: \colorhref{https://github.com/matze/mtheme}
        {https://github.com/matze/mtheme}.

        Projekt typograficzny: Iwona Grabska-Gradzińska;
        Skład: Kamil Ziemian \\
        Korekta: Wojciech Palacz;
        Modele: Dariusz Frymus, Kamil Nowakowski \\
        Rysunki i~wykresy: Kamil Ziemian, Paweł Węgrzyn, Wojciech Palacz \\
        Montaż: Agencja Filmowa Film \& Television Production~-- Zbigniew
        Masklak

      \end{flushleft}

    \end{textblock}

  \end{frame}





  \begin{frame}[standout]


    \begingroup

    \color{jFrametitleFGColor}

    #1

    \endgroup

  \end{frame}
}










% ---------------------------------------
% Commands for lectures "Geometria 3D dla twórców gier wideo"
% English version
% ---------------------------------------
% \newcommand{\FundingLogoWhitePictureEN}
% {./PresentationPictures/CommonPictures/logotypFundusze_biale_bez_tla2.pdf}
\newcommand{\FundingLogoColorPictureEN}
{./PresentationPictures/CommonPictures/European_Funds_color_EN.pdf}
% \newcommand{\EULogoWhitePictureEN}
% {./PresentationPictures/CommonPictures/logotypUE_biale_bez_tla2.pdf}
\newcommand{\EUSocialFundLogoColorPictureEN}
{./PresentationPictures/CommonPictures/EU_Social_Fund_color_EN.pdf}
% \newcommand{\ZintegrUJLogoWhitePictureEN}
% {./PresentationPictures/CommonPictures/zintegruj-logo-white.pdf}
\newcommand{\ZintegrUJLogoColorPictureEN}
{./PresentationPictures/CommonPictures/ZintegrUJ_color.pdf}
\newcommand{\JULogoColorPictureEN}
{./JagiellonianPictures/LogoJU_EN/LogoJU_A_color.pdf}



\newcommand{\GeometryThreeDSpecialBeginningSlideEN}{%
  \begin{frame}[standout]

    \begin{textblock}{11}(1,0.7)

      \begin{flushleft}

        \mdseries

        \footnotesize

        \color{jFrametitleFGColor}

        This content was created as part of a project co-financed by the
        European Union within the framework of the European Social Fund
        POWR.03.05.00-00-Z309/17-00.

      \end{flushleft}

    \end{textblock}





    \begin{textblock}{10}(0,2.2)

      \tikz \fill[color=jBackgroundStyleLight] (0,0) rectangle (12.8,-1.5);

    \end{textblock}


    \begin{textblock}{3.2}(0.7,2.45)

      \includegraphics[scale=0.3]{\FundingLogoColorPictureEN}

    \end{textblock}


    \begin{textblock}{2.5}(4.15,2.5)

      \includegraphics[scale=0.2]{\JULogoColorPictureEN}

    \end{textblock}


    \begin{textblock}{2.5}(6.35,2.4)

      \includegraphics[scale=0.1]{\ZintegrUJLogoColorPictureEN}

    \end{textblock}


    \begin{textblock}{4.2}(8.4,2.6)

      \includegraphics[scale=0.3]{\EUSocialFundLogoColorPictureEN}

    \end{textblock}

  \end{frame}
}



\newcommand{\GeometryThreeDTwoSpecialBeginningSlidesEN}{%
  \begin{frame}[standout]

    \begin{textblock}{11}(1,0.7)

      \begin{flushleft}

        \mdseries

        \footnotesize

        \color{jFrametitleFGColor}

        This content was created as part of a project co-financed by the
        European Union within the framework of the European Social Fund
        POWR.03.05.00-00-Z309/17-00.

      \end{flushleft}

    \end{textblock}





    \begin{textblock}{10}(0,2.2)

      \tikz \fill[color=jBackgroundStyleLight] (0,0) rectangle (12.8,-1.5);

    \end{textblock}


    \begin{textblock}{3.2}(0.7,2.45)

      \includegraphics[scale=0.3]{\FundingLogoColorPictureEN}

    \end{textblock}


    \begin{textblock}{2.5}(4.15,2.5)

      \includegraphics[scale=0.2]{\JULogoColorPictureEN}

    \end{textblock}


    \begin{textblock}{2.5}(6.35,2.4)

      \includegraphics[scale=0.1]{\ZintegrUJLogoColorPictureEN}

    \end{textblock}


    \begin{textblock}{4.2}(8.4,2.6)

      \includegraphics[scale=0.3]{\EUSocialFundLogoColorPictureEN}

    \end{textblock}

  \end{frame}





  \TitleSlideWithPicture
}



\newcommand{\GeometryThreeDSpecialEndingSlideEN}{%
  \begin{frame}[standout]

    \begin{textblock}{11}(1,0.7)

      \begin{flushleft}

        \mdseries

        \footnotesize

        \color{jFrametitleFGColor}

        This content was created as part of a project co-financed by the
        European Union within the framework of the European Social Fund
        POWR.03.05.00-00-Z309/17-00.

      \end{flushleft}

    \end{textblock}





    \begin{textblock}{10}(0,2.2)

      \tikz \fill[color=jBackgroundStyleLight] (0,0) rectangle (12.8,-1.5);

    \end{textblock}


    \begin{textblock}{3.2}(0.7,2.45)

      \includegraphics[scale=0.3]{\FundingLogoColorPictureEN}

    \end{textblock}


    \begin{textblock}{2.5}(4.15,2.5)

      \includegraphics[scale=0.2]{\JULogoColorPictureEN}

    \end{textblock}


    \begin{textblock}{2.5}(6.35,2.4)

      \includegraphics[scale=0.1]{\ZintegrUJLogoColorPictureEN}

    \end{textblock}


    \begin{textblock}{4.2}(8.4,2.6)

      \includegraphics[scale=0.3]{\EUSocialFundLogoColorPictureEN}

    \end{textblock}





    \begin{textblock}{11}(1,4)

      \begin{flushleft}

        \mdseries

        \footnotesize

        \RaggedRight

        \color{jFrametitleFGColor}

        The content of this lecture is made available under a~Creative
        Commons licence (\textsc{cc}), giving the author the credits
        (\textsc{by}) and putting an obligation to share on the same terms
        (\textsc{sa}). Figures and diagrams included in the lecture are
        authored by Paweł Węgrzyn et~al., and are available under the same
        license unless indicated otherwise.\\ The presentation uses the
        Beamer Jagiellonian theme based on Matthias Vogelgesang’s
        Metropolis theme, available under license \LaTeX{} Project
        Public License~1.3c at: \colorhref{https://github.com/matze/mtheme}
        {https://github.com/matze/mtheme}.

        Typographic design: Iwona Grabska-Gradzińska \\
        \LaTeX{} Typesetting: Kamil Ziemian \\
        Proofreading: Wojciech Palacz,
        Monika Stawicka \\
        3D Models: Dariusz Frymus, Kamil Nowakowski \\
        Figures and charts: Kamil Ziemian, Paweł Węgrzyn, Wojciech Palacz

      \end{flushleft}

    \end{textblock}

  \end{frame}
}



\newcommand{\GeometryThreeDTwoSpecialEndingSlidesEN}[1]{%
  \begin{frame}[standout]


    \begin{textblock}{11}(1,0.7)

      \begin{flushleft}

        \mdseries

        \footnotesize

        \color{jFrametitleFGColor}

        This content was created as part of a project co-financed by the
        European Union within the framework of the European Social Fund
        POWR.03.05.00-00-Z309/17-00.

      \end{flushleft}

    \end{textblock}





    \begin{textblock}{10}(0,2.2)

      \tikz \fill[color=jBackgroundStyleLight] (0,0) rectangle (12.8,-1.5);

    \end{textblock}


    \begin{textblock}{3.2}(0.7,2.45)

      \includegraphics[scale=0.3]{\FundingLogoColorPictureEN}

    \end{textblock}


    \begin{textblock}{2.5}(4.15,2.5)

      \includegraphics[scale=0.2]{\JULogoColorPictureEN}

    \end{textblock}


    \begin{textblock}{2.5}(6.35,2.4)

      \includegraphics[scale=0.1]{\ZintegrUJLogoColorPictureEN}

    \end{textblock}


    \begin{textblock}{4.2}(8.4,2.6)

      \includegraphics[scale=0.3]{\EUSocialFundLogoColorPictureEN}

    \end{textblock}





    \begin{textblock}{11}(1,4)

      \begin{flushleft}

        \mdseries

        \footnotesize

        \RaggedRight

        \color{jFrametitleFGColor}

        The content of this lecture is made available under a~Creative
        Commons licence (\textsc{cc}), giving the author the credits
        (\textsc{by}) and putting an obligation to share on the same terms
        (\textsc{sa}). Figures and diagrams included in the lecture are
        authored by Paweł Węgrzyn et~al., and are available under the same
        license unless indicated otherwise.\\ The presentation uses the
        Beamer Jagiellonian theme based on Matthias Vogelgesang’s
        Metropolis theme, available under license \LaTeX{} Project
        Public License~1.3c at: \colorhref{https://github.com/matze/mtheme}
        {https://github.com/matze/mtheme}.

        Typographic design: Iwona Grabska-Gradzińska \\
        \LaTeX{} Typesetting: Kamil Ziemian \\
        Proofreading: Wojciech Palacz,
        Monika Stawicka \\
        3D Models: Dariusz Frymus, Kamil Nowakowski \\
        Figures and charts: Kamil Ziemian, Paweł Węgrzyn, Wojciech Palacz

      \end{flushleft}

    \end{textblock}

  \end{frame}





  \begin{frame}[standout]

    \begingroup

    \color{jFrametitleFGColor}

    #1

    \endgroup

  \end{frame}
}



\newcommand{\GeometryThreeDSpecialEndingSlideVideoVerOneEN}{%
  \begin{frame}[standout]

    \begin{textblock}{11}(1,0.7)

      \begin{flushleft}

        \mdseries

        \footnotesize

        \color{jFrametitleFGColor}

        This content was created as part of a project co-financed by the
        European Union within the framework of the European Social Fund
        POWR.03.05.00-00-Z309/17-00.

      \end{flushleft}

    \end{textblock}





    \begin{textblock}{10}(0,2.2)

      \tikz \fill[color=jBackgroundStyleLight] (0,0) rectangle (12.8,-1.5);

    \end{textblock}


    \begin{textblock}{3.2}(0.7,2.45)

      \includegraphics[scale=0.3]{\FundingLogoColorPictureEN}

    \end{textblock}


    \begin{textblock}{2.5}(4.15,2.5)

      \includegraphics[scale=0.2]{\JULogoColorPictureEN}

    \end{textblock}


    \begin{textblock}{2.5}(6.35,2.4)

      \includegraphics[scale=0.1]{\ZintegrUJLogoColorPictureEN}

    \end{textblock}


    \begin{textblock}{4.2}(8.4,2.6)

      \includegraphics[scale=0.3]{\EUSocialFundLogoColorPictureEN}

    \end{textblock}





    \begin{textblock}{11}(1,4)

      \begin{flushleft}

        \mdseries

        \footnotesize

        \RaggedRight

        \color{jFrametitleFGColor}

        The content of this lecture is made available under a Creative
        Commons licence (\textsc{cc}), giving the author the credits
        (\textsc{by}) and putting an obligation to share on the same terms
        (\textsc{sa}). Figures and diagrams included in the lecture are
        authored by Paweł Węgrzyn et~al., and are available under the same
        license unless indicated otherwise.\\ The presentation uses the
        Beamer Jagiellonian theme based on Matthias Vogelgesang’s
        Metropolis theme, available under license \LaTeX{} Project
        Public License~1.3c at: \colorhref{https://github.com/matze/mtheme}
        {https://github.com/matze/mtheme}.

        Typographic design: Iwona Grabska-Gradzińska;
        \LaTeX{} Typesetting: Kamil Ziemian \\
        Proofreading: Wojciech Palacz,
        Monika Stawicka \\
        3D Models: Dariusz Frymus, Kamil Nowakowski \\
        Figures and charts: Kamil Ziemian, Paweł Węgrzyn, Wojciech
        Palacz \\
        Film editing: Agencja Filmowa Film \& Television Production~--
        Zbigniew Masklak

      \end{flushleft}

    \end{textblock}

  \end{frame}
}



\newcommand{\GeometryThreeDSpecialEndingSlideVideoVerTwoEN}{%
  \begin{frame}[standout]

    \begin{textblock}{11}(1,0.7)

      \begin{flushleft}

        \mdseries

        \footnotesize

        \color{jFrametitleFGColor}

        This content was created as part of a project co-financed by the
        European Union within the framework of the European Social Fund
        POWR.03.05.00-00-Z309/17-00.

      \end{flushleft}

    \end{textblock}





    \begin{textblock}{10}(0,2.2)

      \tikz \fill[color=jBackgroundStyleLight] (0,0) rectangle (12.8,-1.5);

    \end{textblock}


    \begin{textblock}{3.2}(0.7,2.45)

      \includegraphics[scale=0.3]{\FundingLogoColorPictureEN}

    \end{textblock}


    \begin{textblock}{2.5}(4.15,2.5)

      \includegraphics[scale=0.2]{\JULogoColorPictureEN}

    \end{textblock}


    \begin{textblock}{2.5}(6.35,2.4)

      \includegraphics[scale=0.1]{\ZintegrUJLogoColorPictureEN}

    \end{textblock}


    \begin{textblock}{4.2}(8.4,2.6)

      \includegraphics[scale=0.3]{\EUSocialFundLogoColorPictureEN}

    \end{textblock}





    \begin{textblock}{11}(1,4)

      \begin{flushleft}

        \mdseries

        \footnotesize

        \RaggedRight

        \color{jFrametitleFGColor}

        The content of this lecture is made available under a Creative
        Commons licence (\textsc{cc}), giving the author the credits
        (\textsc{by}) and putting an obligation to share on the same terms
        (\textsc{sa}). Figures and diagrams included in the lecture are
        authored by Paweł Węgrzyn et~al., and are available under the same
        license unless indicated otherwise.\\ The presentation uses the
        Beamer Jagiellonian theme based on Matthias Vogelgesang’s
        Metropolis theme, available under license \LaTeX{} Project
        Public License~1.3c at: \colorhref{https://github.com/matze/mtheme}
        {https://github.com/matze/mtheme}.

        Typographic design: Iwona Grabska-Gradzińska;
        \LaTeX{} Typesetting: Kamil Ziemian \\
        Proofreading: Wojciech Palacz,
        Monika Stawicka \\
        3D Models: Dariusz Frymus, Kamil Nowakowski \\
        Figures and charts: Kamil Ziemian, Paweł Węgrzyn, Wojciech
        Palacz \\
        Film editing: IMAVI -- Joanna Kozakiewicz, Krzysztof Magda, Nikodem
        Frodyma

      \end{flushleft}

    \end{textblock}

  \end{frame}
}



\newcommand{\GeometryThreeDSpecialEndingSlideVideoVerThreeEN}{%
  \begin{frame}[standout]

    \begin{textblock}{11}(1,0.7)

      \begin{flushleft}

        \mdseries

        \footnotesize

        \color{jFrametitleFGColor}

        This content was created as part of a project co-financed by the
        European Union within the framework of the European Social Fund
        POWR.03.05.00-00-Z309/17-00.

      \end{flushleft}

    \end{textblock}





    \begin{textblock}{10}(0,2.2)

      \tikz \fill[color=jBackgroundStyleLight] (0,0) rectangle (12.8,-1.5);

    \end{textblock}


    \begin{textblock}{3.2}(0.7,2.45)

      \includegraphics[scale=0.3]{\FundingLogoColorPictureEN}

    \end{textblock}


    \begin{textblock}{2.5}(4.15,2.5)

      \includegraphics[scale=0.2]{\JULogoColorPictureEN}

    \end{textblock}


    \begin{textblock}{2.5}(6.35,2.4)

      \includegraphics[scale=0.1]{\ZintegrUJLogoColorPictureEN}

    \end{textblock}


    \begin{textblock}{4.2}(8.4,2.6)

      \includegraphics[scale=0.3]{\EUSocialFundLogoColorPictureEN}

    \end{textblock}





    \begin{textblock}{11}(1,4)

      \begin{flushleft}

        \mdseries

        \footnotesize

        \RaggedRight

        \color{jFrametitleFGColor}

        The content of this lecture is made available under a Creative
        Commons licence (\textsc{cc}), giving the author the credits
        (\textsc{by}) and putting an obligation to share on the same terms
        (\textsc{sa}). Figures and diagrams included in the lecture are
        authored by Paweł Węgrzyn et~al., and are available under the same
        license unless indicated otherwise.\\ The presentation uses the
        Beamer Jagiellonian theme based on Matthias Vogelgesang’s
        Metropolis theme, available under license \LaTeX{} Project
        Public License~1.3c at: \colorhref{https://github.com/matze/mtheme}
        {https://github.com/matze/mtheme}.

        Typographic design: Iwona Grabska-Gradzińska;
        \LaTeX{} Typesetting: Kamil Ziemian \\
        Proofreading: Wojciech Palacz,
        Monika Stawicka \\
        3D Models: Dariusz Frymus, Kamil Nowakowski \\
        Figures and charts: Kamil Ziemian, Paweł Węgrzyn, Wojciech
        Palacz \\
        Film editing: Agencja Filmowa Film \& Television Production~--
        Zbigniew Masklak \\
        Film editing: IMAVI -- Joanna Kozakiewicz, Krzysztof Magda, Nikodem
        Frodyma

      \end{flushleft}

    \end{textblock}

  \end{frame}
}



\newcommand{\GeometryThreeDTwoSpecialEndingSlidesVideoVerOneEN}[1]{%
  \begin{frame}[standout]

    \begin{textblock}{11}(1,0.7)

      \begin{flushleft}

        \mdseries

        \footnotesize

        \color{jFrametitleFGColor}

        This content was created as part of a project co-financed by the
        European Union within the framework of the European Social Fund
        POWR.03.05.00-00-Z309/17-00.

      \end{flushleft}

    \end{textblock}





    \begin{textblock}{10}(0,2.2)

      \tikz \fill[color=jBackgroundStyleLight] (0,0) rectangle (12.8,-1.5);

    \end{textblock}


    \begin{textblock}{3.2}(0.7,2.45)

      \includegraphics[scale=0.3]{\FundingLogoColorPictureEN}

    \end{textblock}


    \begin{textblock}{2.5}(4.15,2.5)

      \includegraphics[scale=0.2]{\JULogoColorPictureEN}

    \end{textblock}


    \begin{textblock}{2.5}(6.35,2.4)

      \includegraphics[scale=0.1]{\ZintegrUJLogoColorPictureEN}

    \end{textblock}


    \begin{textblock}{4.2}(8.4,2.6)

      \includegraphics[scale=0.3]{\EUSocialFundLogoColorPictureEN}

    \end{textblock}





    \begin{textblock}{11}(1,4)

      \begin{flushleft}

        \mdseries

        \footnotesize

        \RaggedRight

        \color{jFrametitleFGColor}

        The content of this lecture is made available under a Creative
        Commons licence (\textsc{cc}), giving the author the credits
        (\textsc{by}) and putting an obligation to share on the same terms
        (\textsc{sa}). Figures and diagrams included in the lecture are
        authored by Paweł Węgrzyn et~al., and are available under the same
        license unless indicated otherwise.\\ The presentation uses the
        Beamer Jagiellonian theme based on Matthias Vogelgesang’s
        Metropolis theme, available under license \LaTeX{} Project
        Public License~1.3c at: \colorhref{https://github.com/matze/mtheme}
        {https://github.com/matze/mtheme}.

        Typographic design: Iwona Grabska-Gradzińska;
        \LaTeX{} Typesetting: Kamil Ziemian \\
        Proofreading: Wojciech Palacz,
        Monika Stawicka \\
        3D Models: Dariusz Frymus, Kamil Nowakowski \\
        Figures and charts: Kamil Ziemian, Paweł Węgrzyn,
        Wojciech Palacz \\
        Film editing: Agencja Filmowa Film \& Television Production~--
        Zbigniew Masklak

      \end{flushleft}

    \end{textblock}

  \end{frame}





  \begin{frame}[standout]


    \begingroup

    \color{jFrametitleFGColor}

    #1

    \endgroup

  \end{frame}
}



\newcommand{\GeometryThreeDTwoSpecialEndingSlidesVideoVerTwoEN}[1]{%
  \begin{frame}[standout]

    \begin{textblock}{11}(1,0.7)

      \begin{flushleft}

        \mdseries

        \footnotesize

        \color{jFrametitleFGColor}

        This content was created as part of a project co-financed by the
        European Union within the framework of the European Social Fund
        POWR.03.05.00-00-Z309/17-00.

      \end{flushleft}

    \end{textblock}





    \begin{textblock}{10}(0,2.2)

      \tikz \fill[color=jBackgroundStyleLight] (0,0) rectangle (12.8,-1.5);

    \end{textblock}


    \begin{textblock}{3.2}(0.7,2.45)

      \includegraphics[scale=0.3]{\FundingLogoColorPictureEN}

    \end{textblock}


    \begin{textblock}{2.5}(4.15,2.5)

      \includegraphics[scale=0.2]{\JULogoColorPictureEN}

    \end{textblock}


    \begin{textblock}{2.5}(6.35,2.4)

      \includegraphics[scale=0.1]{\ZintegrUJLogoColorPictureEN}

    \end{textblock}


    \begin{textblock}{4.2}(8.4,2.6)

      \includegraphics[scale=0.3]{\EUSocialFundLogoColorPictureEN}

    \end{textblock}





    \begin{textblock}{11}(1,4)

      \begin{flushleft}

        \mdseries

        \footnotesize

        \RaggedRight

        \color{jFrametitleFGColor}

        The content of this lecture is made available under a Creative
        Commons licence (\textsc{cc}), giving the author the credits
        (\textsc{by}) and putting an obligation to share on the same terms
        (\textsc{sa}). Figures and diagrams included in the lecture are
        authored by Paweł Węgrzyn et~al., and are available under the same
        license unless indicated otherwise.\\ The presentation uses the
        Beamer Jagiellonian theme based on Matthias Vogelgesang’s
        Metropolis theme, available under license \LaTeX{} Project
        Public License~1.3c at: \colorhref{https://github.com/matze/mtheme}
        {https://github.com/matze/mtheme}.

        Typographic design: Iwona Grabska-Gradzińska;
        \LaTeX{} Typesetting: Kamil Ziemian \\
        Proofreading: Wojciech Palacz,
        Monika Stawicka \\
        3D Models: Dariusz Frymus, Kamil Nowakowski \\
        Figures and charts: Kamil Ziemian, Paweł Węgrzyn,
        Wojciech Palacz \\
        Film editing: IMAVI -- Joanna Kozakiewicz, Krzysztof Magda, Nikodem
        Frodyma

      \end{flushleft}

    \end{textblock}

  \end{frame}





  \begin{frame}[standout]


    \begingroup

    \color{jFrametitleFGColor}

    #1

    \endgroup

  \end{frame}
}



\newcommand{\GeometryThreeDTwoSpecialEndingSlidesVideoVerThreeEN}[1]{%
  \begin{frame}[standout]

    \begin{textblock}{11}(1,0.7)

      \begin{flushleft}

        \mdseries

        \footnotesize

        \color{jFrametitleFGColor}

        This content was created as part of a project co-financed by the
        European Union within the framework of the European Social Fund
        POWR.03.05.00-00-Z309/17-00.

      \end{flushleft}

    \end{textblock}





    \begin{textblock}{10}(0,2.2)

      \tikz \fill[color=jBackgroundStyleLight] (0,0) rectangle (12.8,-1.5);

    \end{textblock}


    \begin{textblock}{3.2}(0.7,2.45)

      \includegraphics[scale=0.3]{\FundingLogoColorPictureEN}

    \end{textblock}


    \begin{textblock}{2.5}(4.15,2.5)

      \includegraphics[scale=0.2]{\JULogoColorPictureEN}

    \end{textblock}


    \begin{textblock}{2.5}(6.35,2.4)

      \includegraphics[scale=0.1]{\ZintegrUJLogoColorPictureEN}

    \end{textblock}


    \begin{textblock}{4.2}(8.4,2.6)

      \includegraphics[scale=0.3]{\EUSocialFundLogoColorPictureEN}

    \end{textblock}





    \begin{textblock}{11}(1,4)

      \begin{flushleft}

        \mdseries

        \footnotesize

        \RaggedRight

        \color{jFrametitleFGColor}

        The content of this lecture is made available under a Creative
        Commons licence (\textsc{cc}), giving the author the credits
        (\textsc{by}) and putting an obligation to share on the same terms
        (\textsc{sa}). Figures and diagrams included in the lecture are
        authored by Paweł Węgrzyn et~al., and are available under the same
        license unless indicated otherwise. \\ The presentation uses the
        Beamer Jagiellonian theme based on Matthias Vogelgesang’s
        Metropolis theme, available under license \LaTeX{} Project
        Public License~1.3c at: \colorhref{https://github.com/matze/mtheme}
        {https://github.com/matze/mtheme}.

        Typographic design: Iwona Grabska-Gradzińska;
        \LaTeX{} Typesetting: Kamil Ziemian \\
        Proofreading: Leszek Hadasz, Wojciech Palacz,
        Monika Stawicka \\
        3D Models: Dariusz Frymus, Kamil Nowakowski \\
        Figures and charts: Kamil Ziemian, Paweł Węgrzyn,
        Wojciech Palacz \\
        Film editing: Agencja Filmowa Film \& Television Production~--
        Zbigniew Masklak \\
        Film editing: IMAVI -- Joanna Kozakiewicz, Krzysztof Magda, Nikodem
        Frodyma


      \end{flushleft}

    \end{textblock}

  \end{frame}





  \begin{frame}[standout]


    \begingroup

    \color{jFrametitleFGColor}

    #1

    \endgroup

  \end{frame}
}











% ------------------------------------------------------
% BibLaTeX
% ------------------------------------------------------
% Package biblatex, with biber as its backend, allow us to handle
% bibliography entries that use Unicode symbols outside ASCII.
\usepackage[
language=polish,
backend=biber,
style=alphabetic,
url=false,
eprint=true,
]{biblatex}

\addbibresource{Podstawy-informatyki-ETC-Bibliography.bib}





% ------------------------------------------------------
% Importing packages, libraries and setting their configuration
% ------------------------------------------------------
% Library improving positioning of nodes in graphs
% \usetikzlibrary{positioning}





% ------------------------------------------------------
% Local packages
% ------------------------------------------------------
% Local configuration of this particular presentation.
\usepackage{./Local-packages/local-settings}

% Package containing various command useful for working with a text.
\usepackage{./Local-packages/general-commands}

% % Styles for drawing diagrams
% \usepackage{./Local-packages/PGF-TikZ-Diagram-styles}

% % Jagiellonian theme's colors
% \usepackage{./Local-packages/jagiellonian-theme-colors}



% ######################################
% Local macros
% ######################################



% Macros for writing numbers
\newcommand{\btext}{\text{b}}
\newcommand{\otext}{\text{o}}
\newcommand{\xtext}{\text{x}}

\newcommand{\Atext}{\text{A}}
\newcommand{\Btext}{\text{B}}
\newcommand{\Ctext}{\text{C}}
\newcommand{\Dtext}{\text{D}}
\newcommand{\Etext}{\text{E}}
\newcommand{\Ftext}{\text{F}}

\newcommand{\horSpaceOneLoc}{\hspace{3em}}










% ####################################################################
\title{Podstawy informatyki z~językiem~C}
\subtitle{Programowanie nie jest trudne, debugowanie już tak}

\author{Kamil Ziemian \\
  \email}


% \date{}
% ####################################################################










% ####################################################################
% Beginning of the document
\begin{document}
% ####################################################################





% ######################################
% Number of chars: 7k+, 29k+,
% Text is adjusted to the left and words are broken at the end of the line.
\RaggedRight
% ######################################





% ######################################
\maketitle
% ######################################





% ##################
\begin{frame}
  \frametitle{Spis treści}


  \tableofcontents

\end{frame}
% ##################










% ################################################
\section{O~liczbach w~komputerze}
% ################################################


% ######################################
\begin{frame}
  \frametitle{Komputer to maszyna do liczenia}


  W~samej angielskiej nazwie \textit{computer} zawiera~się informacja,
  że~jest to maszyna przeznaczona do liczenia, więc siłą rzeczy należy
  coś powiedzieć o~tym jak komputer zarządza liczbami. Nie będziemy przy
  tym omawiali wszystkich zawiłości jakie dotyczą tego problemu. Na tym
  przedmiocie chcemy przekazać tylko tyle informacji, ile jest
  naszym zdaniem potrzebne do wyrobienia sobie podstawowych intuicji.
  Za głębsze wniknięcie w~naturę rzeczy odpowiedzialne są inne przedmioty.

  Jednym z~powodów dla których zaczynamy od liczb, jest to, że~gdy omówimy
  jak komputer przechowuje liczby, będziemy mogli zrozumieć, jak przechowuje
  inne rzeczy. Dzieje~się tak, że~dla komputera wszystko jest ciągiem
  bajtów, a~każdy bajt można zrozumieć jako liczbę całkowitą, z~zakresu
  $0$, $1$, $2$, \ldots, $253$, $254$, $255$. Może~się to wydawać trochę
  nieintuicyjne, proszę jednak uwierzyć, że~to ma sens i~dać nam trochę
  czasu na wytłumaczenie tego.

\end{frame}
% ######################################





% ######################################
\begin{frame}
  \frametitle{System pozycyjny}


  Najprostsze liczby to liczby naturalne: $0$, $1$, $2$, \ldots, dlatego
  naszą dyskusję zaczniemy właśnie od nich. Niektórzy nie zaliczają
  liczby~$0$ do liczb naturalnych, jednak ze~względu na użyteczność tej
  liczby w~obliczeniach, w~informatyce przyjmujemy konwencję, że~$0$ jest
  liczbą naturalną.

  Komputery pracują w~systemie binarnym, co wynika zapewne z~tego powodu,
  że~komputer są urządzeniami elektronicznymi, a~choć można zbudować
  urządzenia które pracują choćby w~systemie dziesiętnym, to jest to dużo
  trudniejsze. To jednak temat na zajęcia z~elektroniki.

  Niemniej, warto zacząć od przyjrzenia~się dobrze nam znanemu systemu
  dziesiętnemu, gdyż do niego jesteśmy najbardziej przyzwyczajeniu.
  Jest on systemem pozycyjnym zapisu liczb o~podstawie równej liczbie
  dziesięć. W~systemie tym mamy dziesięć bazowych symboli, zwanych cyframi:
  $0$, $1$, $2$, $3$, $4$, $5$, $6$, $7$, $8$ i~$9$.

\end{frame}
% ######################################





% ######################################
\begin{frame}
  \frametitle{System pozycyjny}


  Za pomocą tego ciągu cyfr możemy zapisać dowolną liczbę naturalną. Przy
  tym proszę zwrócić uwagę, na to, iż~sama liczba dziesięć \alert{nie jest}
  reprezentowana przez żadną cyfrę, lecz przez ciąg dwóch cyfr~$10$.
  Warto o~tym pamiętać, porównując ten system z~systemem binarnym.

  Ten system zapisu liczb nosi nazwę systemu pozycyjnego, bo liczby
  naturalne zapisuje jako ciągi cyfr, a~następnie czytamy od prawej
  do~lewej, przy czym znaczenie każdej cyfry zależy od~jej \alert{pozycji}
  w~tym ciągu. By to lepiej zobaczyć, rozpatrzmy kilka przykładów.

  Zanim podamy przykłady, należy wspomnieć, że~dla zgodności z~notacją
  z~języka~C, znak równa~się zapisujemy jako $==$. Przypominamy,
  że~w~języku C symbol „$=$” oznacza \alert{przypisanie} wartości
  do~zmiennej, nie zaś równość dwóch liczb.

\end{frame}
% ######################################





% ######################################
\begin{frame}
  \frametitle{System pozycyjny}


  Kilka przykładów czytania liczb w~systemie dziesiętnym. \\
  $123 == 3 \cdot 10^{ \HorSpaceTwo 0 } + 2 \cdot 10^{ \HorSpaceTwo 1 } +
  1 \cdot 10^{ \HorSpaceTwo 2 } == 3 \cdot 1 + 2 \cdot 10 + 1 \cdot 100 == 3 + 20 + 100$. \\
  $456 == 6 \cdot 10^{ \HorSpaceTwo 0 } + 5 \cdot 10^{ \HorSpaceTwo 1 } +
  4 \cdot 10^{ \HorSpaceTwo 2 } == 6 \cdot 1 + 5 \cdot 10 + 4 \cdot 100 == 6 + 50 + 400$. \\
  $2 \, 345 == 5 \cdot 10^{ \HorSpaceTwo 0 } + 4 \cdot 10^{ \HorSpaceTwo 1 } +
  3 \cdot 10^{ \HorSpaceTwo 2 } + 2 \cdot 10^{ \HorSpaceTwo 3 } == 5 \cdot 1 + 4 \cdot 10 +
  3 \cdot 100 \, +$ \\
  $+ \, 2 \cdot 1 \, 000 == 5 + 40 + 300 + 2 \, 000$.

  Dopisywanie zer po lewej stronie danej liczby, nie zmienia jej wartość.
  Czyli $12 == 012 == 0012 == 00012$. W~informatyce ten drobny szczegół
  ma duże znaczenie.

  W~systemie binarnym mamy tylko dwie cyfry $0$ i~$1$. Element pamięci
  komputer, który jest w~stanie przechować jedną z~wartości $0$ lub $1$,
  nazywamy \textbf{bitem}. Wszystko co komputer przechowuje w~pamięci,
  teksty, obrazy, dźwięki, jest dla niego ciągiem bitów. Jak to działa,
  o~tym spróbujemy powiedzieć później.

\end{frame}
% ######################################





% ######################################
\begin{frame}
  \frametitle{Wybrane potęgi liczby~$2$}


  W~systemie dziesiętnym lubimy potęgi liczby dziesięć, które z~tego powodu
  mają w~większości przypadków specjalne nazwy, takie jak „tysiąc”
  czy~„milion”. Dokładnie z~tego samego powodu w~systemie binarnym lubimy
  potęgi liczby dwa. Choć jest to strasznie nudne dla normalnego człowieka,
  przyjrzymy~się teraz pewnym potęgom liczby dwa. Wbrew pozorom,
  w~informatyce ta wiedza potrafi być cenna.

  $2^{ \HorSpaceTwo 0 } == 1$ \\[0.1em]
  $2^{ \HorSpaceTwo 1 } == 2$ \\[0.1em]
  $2^{ \HorSpaceTwo 2 } == 4$ \\[0.1em]
  $2^{ \HorSpaceTwo 3 } == 8$ \\[0.1em]
  $2^{ \HorSpaceTwo 4 } == 16$ \\[0.1em]
  $2^{ \HorSpaceTwo 5 } == 32$ \\[0.1em]
  $2^{ \HorSpaceTwo 6 } == 64$ \\[0.1em]
  $2^{ \HorSpaceTwo 7 } == 128$ \\[0.1em]
  $2^{ \HorSpaceTwo 8 } == 256$

\end{frame}
% ######################################





% ######################################
\begin{frame}
  \frametitle{Wybrane potęgi liczby~$2$}


  Podstawową jednostką pamięci jaką dziś używamy w~komputerach, jest
  jeden bajt, czyli $8 == 2^{ 3 }$ bitów. To jeden z~przykładów tego jak
  lubimy potęgi liczby~$2$.

  $2^{ \HorSpaceTwo 9 } == 512$ \\[0.1em]
  $2^{ \HorSpaceTwo 10 } == 1 \; 024$ \\[0.1em]
  $2^{ \HorSpaceTwo 11 } == 2 \; 048$ \\[0.1em]
  $2^{ \HorSpaceTwo 12 } == 4 \; 096$ \\[0.1em]
  $2^{ \HorSpaceTwo 15 } == 32 \; 768$ \\[0.1em]
  $2^{ \HorSpaceTwo 16 } == 65 \; 536$

  Liczby $2^{ \HorSpaceTwo 16 } == 65 \; 536$
  i~$2^{ \HorSpaceTwo 16 } - 1 == 65 \; 535$ pojawiają~się
  często w~różnych działach techniki, gdyż często zachodzi potrzeba
  przesyłania informacji, której podstawową jednostką są dwa bajty,
  czyli $16$~bitów.

\end{frame}
% ######################################





% ######################################
\begin{frame}
  \frametitle{Wybrane potęgi liczby~$2$}


  $2^{ \HorSpaceTwo 17 } == 131 \; 072$ \\[0.1em]
  $2^{ \HorSpaceTwo 18 } == 262 \; 144$ \\[0.1em]
  $2^{ \HorSpaceTwo 19 } == 524 \; 288$ \\[0.1em]
  $2^{ \HorSpaceTwo 20 } == 1 \; 048 \; 576$ \\[0.1em]
  $2^{ \HorSpaceTwo 21 } == 2 \; 097 \; 152$ \\[0.1em]
  $2^{ \HorSpaceTwo 22 } == 4 \; 194 \; 304$ \\[0.1em]
  $2^{ \HorSpaceTwo 23 } == 8 \; 388 \; 608$ \\[0.1em]
  $2^{ \HorSpaceTwo 24 } == 16 \; 777 \; 216$ \\[0.1em]
  $2^{ \HorSpaceTwo 25 } == 33 \; 554 \; 432$ \\[0.1em]
  $2^{ \HorSpaceTwo 26 } == 67 \; 108 \; 864$ \\[0.1em]
  $2^{ \HorSpaceTwo 27 } == 134 \; 217 \; 728$ \\[0.1em]
  $2^{ \HorSpaceTwo 28 } == 268 \; 435 \; 456$ \\[0.1em]
  $2^{ \HorSpaceTwo 29 } == 536 \; 870 \; 912$ \\[0.1em]
  $2^{ \HorSpaceTwo 30 } == 1 \; 073 \; 741 \; 824$ \\[0.1em]
  $2^{ \HorSpaceTwo 31 } == 2 \; 147 \; 483 \; 648$

\end{frame}
% ######################################





% ######################################
\begin{frame}
  \frametitle{Wybrane potęgi liczby~$2$}


  $2^{ \HorSpaceTwo 32 } == 4 \; 294 \; 967 \; 296$

  Liczba $2^{ \HorSpaceTwo 32 } == 4 \; 294 \; 967 \; 296$ jest bardzo ważna
  dla wszystkich \\
  systemów i~programów $32$-bitowych, czyli takich w~których
  naturalną jednostką przesyłu informacji są właśnie $32$ bity ($4$~bajty).

  $2^{ \HorSpaceTwo 63 } == 9 \; 223 \; 372 \; 036 \; 854 \; 775 \; 808 \approx
  9 \cdot 10^{ 18 }$ \\
  Słownie: dziewięć trylionów, dwieście dwadzieścia trzy biliardy, trzysta
  siedemdziesiąt dwa biliony, trzydzieści sześć miliardów, osiemset
  pięćdziesiąt cztery miliony, siedemset siedemdziesiąt pięć tysięcy,
  osiemset osiem. \\
  Prościej: $9$~trylionów, $223$~biliardy, $372$~biliony,
  $36$~miliardy, $854$~miliony, $775$~tysięcy, $808$.

\end{frame}
% ######################################





% ######################################
\begin{frame}
  \frametitle{Wybrane potęgi liczby~$2$}


  $2^{ \HorSpaceTwo 64 } == 18 \; 446 \; 744 \; 073 \; 709 \; 551 \; 616 \approx
  1.8 \cdot 10^{ 19 }$ \\
  Słownie: osiemnaście trylionów, czterysta czterdzieści sześć biliardów,
  siedemset czterdzieści cztery biliony, siedemdziesiąt trzy miliardy,
  siedemset dziewięć milionów, pięćset pięćdziesiąt jeden milionów,
  sześćset szesnaście. \\
  Prościej: $18$ trylionów, $446$ biliardów, $744$ biliony, $73$ miliardy,
  $709$ milionów, $551$ tysięcy, $616$.

  Jak ktoś~się tej liczby nauczy na pamięć, to zasługuje na $5$~dodatkowych
  punktów z~prac domowych.

  Liczba ta jest oczywiście bardzo ważna, dla systemów $64$-bitowych.

\end{frame}
% ######################################





% ######################################
\begin{frame}
  \frametitle{System binarny zapisu liczb naturalnych}


  Przejdziemy teraz do tego, jak zapisywać liczby naturalne w~systemie
  binarnym. Wzorując~się na języku programowania Go, liczby binarne
  będziemy zapisywać jako $0\text{b}1$, $0\text{b}01$, $0\text{b}100$,~etc.
  Tak jak w~systemie dziesiętnym, czytamy je od prawej do lewej.
  Prześledźmy teraz kilka przykładów.

  $0\btext 0 == 0 \cdot 2^{ \HorSpaceTwo 0 } == 0 \cdot 1 == 0$. \\
  $0\btext 1 == 1 \cdot 2^{ \HorSpaceTwo 0 } == 1 \cdot 1 == 1$. \\
  $0\btext 10 == 0 \cdot 2^{ \HorSpaceTwo 0 } + 1 \cdot 2^{ \HorSpaceTwo 1 } ==
  0 \cdot 1 + 1 \cdot 2 == 2$. \\
  $0\btext 11 == 1 \cdot 2^{ \HorSpaceTwo 0 } + 1 \cdot 2^{ \HorSpaceTwo 1 } ==
  1 \cdot 1 + 1 \cdot 2 == 1 + 2 == 3$. \\
  $0\btext 10 == 0 \cdot 2^{ \HorSpaceTwo 0 } + 0 \cdot 2^{ \HorSpaceTwo 1 } +
  1 \cdot 2^{ \HorSpaceTwo 2 } == 0 \cdot 1 + 0 \cdot 2 + 1 \cdot 4 == 1 + 2 + 4 == 7$.

  Rozpatrzmy jeszcze kilka innych przykładów. \\
  $0\btext 111 == 1 \cdot 2^{ \HorSpaceTwo 0 } + 1 \cdot 2^{ \HorSpaceTwo 1 } +
  1 \cdot 2^{ \HorSpaceTwo 2 } == 1 \cdot 1 + 1 \cdot 2 + 1 \cdot 4 == 1 + 2 + 4 == 7$. \\
  $0\btext 10101 == 1 \cdot 2^{ \HorSpaceTwo 0 } + 0 \cdot 2^{ \HorSpaceTwo 1 } +
  1 \cdot 2^{ \HorSpaceTwo 2 } + 0 \cdot 2^{ \HorSpaceTwo 3 } + 1 \cdot 2^{ \HorSpaceTwo 4 } ==
  1 \cdot 1 + 0 \cdot 2 + 1 \cdot 4 \, +$ \\
  $+ \, 0 \cdot 8 + 1 \cdot 16 == 1 + 4 + 16 == 21$.

\end{frame}
% ######################################





% ######################################
\begin{frame}
  \frametitle{System binarny zapisu liczb naturalnych}


  Porównajmy teraz system dziesiętny do~binarnego. W~systemie dziesiętnym
  mamy dziesięć cyfr: $0$, $1$, $2$, $3$, $4$, $5$, $6$, $7$, $8$ i~$9$.
  Podstawę tego systemu, liczbę dziesięć zapisujemy jako $10$. W~systemie
  binarnym mamy dwie cyfry: $0$ i~$1$. Podstawę tego systemu, liczbę~$2$
  zapisujemy jako $0\btext 10$.

  Rozpatrzmy teraz ważny dla informatyki problem. Ile liczb możemy zapisać
  za~pomocą jednej cyfry dziesiętnej? Odpowiedź jest, że~dziesięć:
  $0$, $1$, $2$, $3$, $4$, $5$, $6$, $7$, $8$ i~$9$. Ile liczb możemy
  zapisać za pomocą dwóch cyfr dziesiętnych? Liczb takich jest sto:
  $0$, $1$, $2$, \ldots, $97$, $98$, $99$. Analogicznie, za pomocą trzech cyfr
  dziesiętnych możemy zapisać tysiąc liczb: $0$, $1$, $2$, \ldots, $997$,
  $998$, $999$.

  Proszę zauważyć, że~dla jednej cyfry mamy $10^{ \, 1 }$ możliwych liczb,
  dla dwóch $100 == 10^{ \, 2 }$ liczb, dla trzech $1 \, 000 == 10^{ \, 3 }$
  liczb. Ogólnie, za pomocą ciąŋu~$n$ cyfr dziesiętnych możemy zapisać
  $10^{ \, n }$.

\end{frame}
% ######################################





% ######################################
\begin{frame}
  \frametitle{System binarny zapisu liczb naturalnych}


  Ogólniej, jeśli podstawą systemu jest liczba $p$, to za pomocą
  $n$ cyfr tego systemu możemy zapisać $p^{ \, n }$ liczb. Ten matematyczny
  fakt, można prosto udowodnić metodą indukcji, ale nie będziemy tego
  robić, bo dowodzenie taki rzeczy raczej państwa nie interesuje, a~to nie
  są zajęcia z~matematyki.

  Przyjrzyjmy~się jak to działa da systemu binarnego. Dla jednej cyfry
  binarnej (jednego bitu), powinno być $2^{ \, 1 } == 2$ liczb i~tak
  jest w~istocie: $0\btext 0 == 0$ i~$0\btext 1 == 1$. Dla dwóch cyfr
  binarnych (dwóch bitów), powinno być $2^{ \, 2 } == 4$ liczby.
  Oto one: $0\btext 00 == 0$, $0\btext 01 == 1$, $0\btext 10 == 2$
  i~$0\btext 11 == 3$. Dla trzech cyfr binarnych (trzech bitów) mamy
  już $2^{ \, 3 } == 8$ liczb: $0\btext 000 == 0$, $0\btext 001 == 1$,
  $0\btext 010 == 2$, $0\btext 011 == 3$, $0\btext 100 == 4$,
  $0\btext 101 == 5$, $0\btext 110 == 6$ i~$0\btext 111 == 7$.

\end{frame}
% ######################################





% ######################################
\begin{frame}
  \frametitle{System binarny zapisu liczb naturalnych}


  Dla czterech cyfr binarnych (czterech bitów), mamy już $2^{ \, 4 } == 16$
  możliwych liczb. Ten przykład można przeanalizować dokładnie tak samo,
  nie będziemy tego robić tylko ze względu na ilość czasu i~miejsca
  jaki by to zajęło.

  Jak może Państwo już zauważyli, za~pomocą $n$ bitów zapisujemy
  $2^{ \, n }$ liczb: $0$, $1$, \ldots, $2^{ \, n } - 1$. Proszę zwrócić uwagę,
  że~\alert{największą} liczbą jaką w~ten sposób zapisujemy jest
  $2^{ \, n } - 1$.

  Teraz trochę kwestii technicznych. Najmniejszą jednostką pamięci, jakiej
  używa komputer jest jeden bit: $0$ lub~$1$. Jednak ludzkość uznała,
  że~jest bardzo niewygodnie operować na tak małej jednostce pamięci,
  dlatego przyjęto, że~najmniejszą jednostką \alert{przesyłania} danych
  jest \textbf{bajt}. Z~definicji bajt to $8 == 2^{ \, 3 }$ bitów, gdzie
  wychodzi upodobanie informatyki do potęg liczby~$2$.

  Z~tego właśnie powodu mówimy, że~pamięć komputera jest ciągiem bajtów.

\end{frame}
% ######################################





% ######################################
\begin{frame}
  \frametitle{System binarny zapisu liczb naturalnych}


  Zgodnie z~tym co powiedziano wcześniej, każdy bajt może przechowywać
  liczby od $0$ do $2^{ \HorSpaceOne 8 } - 1 == 256 - 1 == 255$. Z~tego
  powodu dowolną rzecz w~komputerze, czy to tekst, czy to obraz, czy
  muzyka, można przedstawić za pomocą ciągu liczby naturalny, przy czym
  każda z~tych liczb jest w~zakresie od~$0$ do~$255$. Bardzo możliwe,
  że~z~tego powodu już wielokrotnie natknęli~się Państwo na te liczby
  w~swoim życiu.

  To może wydawać~się nieprawdopodobne. Jak zakodować powiedzmy utwór
  muzyczny jako ciąg liczb od~$0$ do $255$? Na razie proszę nam uwierzyć,
  że~jest to wykonalne.

  Na razie omówiliśmy tylko jak komputer zapisuje liczby naturalne: $0$,
  $1$, $2$, \ldots A~co z~liczbami ujemnymi? A~z~liczbami wymiernymi?
  A~z~liczbami niewymiernymi? Te temat są bardziej skomplikowane, niż
  mogłoby~się wydawać, więc zanim to zrobimy, omówimy coś prostszego.
  Mianowicie to, jak reprezentować tekst za pomocą liczb.

\end{frame}
% ######################################










% ################################################
\section{Reprezentowanie litera jako liczb}
% ################################################


% ######################################
\begin{frame}
  \frametitle{Reprezentowanie litera jako liczb}


  Zacznijmy od bardzo prostego przykładu, który może~się wydawać
  dziwny, ale jest wbrew pozorom bardzo praktyczny. Jak można zapamiętać
  angielskie słowo „Hello” w~komputerze? Możemy to zrobić w~następujący
  sposób. Wiemy, że~komputer umie zapamiętywać liczby naturalne, więc
  ponumerujmy litery alfabetu (angielskiego) kolejnymi liczbami
  w~następujący sposób. \\
  $\text{A} \to 1$, $\text{B} \to 2$, $\text{C} \to 3$, \ldots, $\text{Z} \to 26$,
  $\text{a} \to 27$, $\text{b} \to 28$, $\text{c} \to 29$, \ldots,
  $\text{z} \to 52$.

  Taka odpowiedniość między literami alfabetu angielskiego, a~liczbami
  naturalnymi nazywamy \textbf{kodowaniem}. W~kontekście tego kursu, przez
  kodowanie będziemy zawsze rozumieli przypisanie pewnym obiektom
  odpowiednich liczb naturalnych.

  Od razu wyjaśnijmy kilka rzeczy. Czemu zajmujemy~się tylko alfabetem
  angielskim? Bo informatykę jaką znamy stworzono w~USA, a~tam głównym
  językiem był i~jest angielski, więc komputerze kodowanie jego alfabetu
  jest wyjątkowo proste.

\end{frame}
% ######################################





% ######################################
\begin{frame}
  \frametitle{Reprezentowanie litera jako liczb}


  Jak kodowane są takie litery jak używane w~języku polski~„ą”, czy symbole
  języka japońskiego, to niestety \alert{dużo} bardziej skomplikowany temat.
  Wrócimy do tego na~krótko potem.

  Alfabet języka angielskiego ma $26$ liter, więc uwzględniając,
  że~mamy duże i~małe litery, w~sumie potrzebujemy liczb od~$1$ do $52$.
  Poniżej można znaleźć tabele opisującą jak działa nasze kodowania
  liczb na litery.

  \begin{center}

    \begin{tabular}{| c | c | c | c | c | c |}
      \hline
      $\text{A} \to 1$ & $\text{B} \to 2$ & $\text{C} \to 3$
      & $\text{D} \to 4$ & $\text{E} \to 5$ & $\text{F} \to 6$ \\
      \hline
      $\text{G} \to 7$ & $\text{H} \to 8$ & $\text{I} \to 9$
      & $\text{J} \to 10$ & $\text{K} \to 11$ & $\text{L} \to 12$ \\
      \hline
      $\text{M} \to 13$ & $\text{N} \to 14$ & $\text{O} \to 15$
      & $\text{P} \to 16$ & $\text{Q} \to 17$ & $\text{R} \to 18$ \\
      \hline
      $\text{S} \to 19$ & $\text{T} \to 20$ & $\text{U} \to 21$
      & $\text{V} \to 22$ & $\text{W} \to 23$ & $\text{X} \to 24$ \\
      \hline
      $\text{Y} \to 25$ & $\text{Z} \to 26$ & & & & \\
      \hline
    \end{tabular}

  \end{center}

\end{frame}
% ######################################





% ######################################
\begin{frame}
  \frametitle{Tabele naszego kodowania}


  \begin{center}

    \begin{tabular}{| c | c | c | c | c | c |}
      \hline
      $\text{a} \to 27$ & $\text{b} \to 28$ & $\text{c} \to 29$
      & $\text{d} \to 30$ & $\text{e} \to 31$ & $\text{f} \to 32$ \\
      \hline
      $\text{g} \to 33$ & $\text{h} \to 34$ & $\text{i} \to 35$
      & $\text{j} \to 36$ & $\text{k} \to 37$ & $\text{l} \to 38$ \\
      \hline
      $\text{m} \to 39$ & $\text{n} \to 40$ & $\text{o} \to 41$
      & $\text{p} \to 42$ & $\text{q} \to 43$ & $\text{r} \to 44$ \\
      \hline
      $\text{s} \to 45$ & $\text{t} \to 46$ & $\text{u} \to 47$
      & $\text{v} \to 48$ & $\text{w} \to 49$ & $\text{x} \to 50$ \\
      \hline
      $\text{y} \to 51$ & $\text{z} \to 52$ & & & & \\
      \hline
    \end{tabular}

  \end{center}

  Możemy teraz zakodować słowo „Hello” jako następujący ciąg liczb:
  $8 \; 31 \; 38 \; 38 \; 41$.

  Jak zobaczymy niedługo, komputer naprawdę
  używa koduje znaki jako liczby, choć nie w~ten sposób, który podaliśmy.
  Dobrze, ale jak komputer wyświetla liczby na ekranie? W~skrócie to działa
  tak, że~komputer dostaje rozkaz „Przełącz~się na tryb, w~którym każdą
  liczbę od $1$ do~$52$ oznacza pewną literę alfabetu angielskiego.
  Jeśli ją napotkasz, to na ekranie wyświetl odpowiedni znak.”. To mamy
  nadzieję, jest wystarczające wyjaśnienie.

\end{frame}
% ######################################





% ######################################
\begin{frame}
  \frametitle{To wciąż za mało}


  Jednak wciąż nie możemy wypisać na ekranie „Hello, World!”, bo umiemy
  tylko kodować liczby, a~co ze znakami przecinka, spacji i~wykrzyknika?
  Nie ma problemu po prostu dorzucamy następujące zasady do naszego kodu,
  gdzie ostatnia pozycja symbolizuje kodowanie znaku spacji.

  \begin{center}

    \begin{tabular}{| c | c | c |}
      \hline
      $\text{,} \to 53$ & $\text{!} \to 54$ & $\text{'\hphantom{a}'} \to 55$ \\
      \hline
    \end{tabular}

  \end{center}

  Teraz możemy zapisać „Hello, World!” jako ciąg liczb: \\
  $8 \; 31 \; 38 \; 38 \; 41 \; 53 \; 23 \; 41 \; 44 \; 38 \; 30 \; 55$

  Świetnie, ale teraz powstaje problem. Co jeśli ktoś uzna, że najpierw
  należy kodować małe litery, a potem duże? Czyli użyje następującego
  kodowania. \\
  $\text{a} \to 1$, $\text{b} \to 2$, $\text{c} \to 3$, \ldots, $\text{z} \to 26$,
  $\text{A} \to 27$, $\text{B} \to 28$, $\text{C} \to 29$, \ldots,
  $\text{Z} \to 52$.

\end{frame}
% ######################################


% a -> 1
% b -> 2
% c -> 3
% d -> 4
% e -> 5
% f -> 6
% g -> 7
% h -> 8
% i -> 9
% j -> 10
% k -> 11
% l -> 12
% m -> 13
% n -> 14
% o -> 15
% p -> 16
% q -> 17
% r -> 18
% s -> 19
% t -> 20
% u -> 21
% v -> 22
% w -> 23
% x -> 24
% y -> 25
% z -> 26
% A -> 27
% B -> 28
% C -> 29
% D -> 30
% E -> 31
% F -> 32
% G -> 33
% H -> 34
% I -> 35
% J -> 36
% K -> 37
% L -> 38
% M -> 39
% N -> 40
% O -> 41
% P -> 42
% Q -> 43
% R -> 44
% S -> 45
% T -> 46
% U -> 47
% V -> 48
% W -> 49
% X -> 50
% Y -> 51
% Z -> 52
% . -> 53
% ! -> 54
%   -> 55





% ######################################
\begin{frame}
  \frametitle{To wciąż za mało}


  A~jeśli ktoś uzna, że~kodowanie powinno być takie, że~najpierw
  duża litera, potem mała? \\
  $\text{A} \to 1$, $\text{a} \to 2$, $\text{B} \to 3$, $\text{b} \to 3$, \ldots,
  $\text{z} \to 26$, $\text{A} \to 27$, $\text{B} \to 28$, $\text{C} \to 29$, \ldots,
  $\text{Z} \to 52$.

\end{frame}
% ######################################




% ##################
\begin{frame}
  \frametitle{Historia kodu ASCII}




  To wbrew pozorom dobrze opisuje, jak komputery działają, co więcej,
  język~C działa tak cały czas. Jednak zanim powiemy sobie jak język~C
  sobie z~tym radzi, to musimy powiedzieć coś o~systemie kodowania
  \colorhref{https://en.wikipedia.org/wiki/ASCII}{\textsc{ascii}}, który
  w~$2025$ wciąż jest jednym z~fundamentów działania naszych komputerów.

  Akronim \colorhref{https://en.wikipedia.org/wiki/ASCII}{\textsc{ascii}}
  pochodzi od~\textit{American Standard Code for~Information Interchange},
  czyli \textit{Standardowy Amerykański Kod do Wymiany Informacji}.
  Jak sama nazwa wskazuje, został on stworzony przez Amerykanów na ich
  potrzeby, co wyjaśnia dlaczego opiera~się na alfabecie angielskim.
  Dlatego też niektórzy uważają, że~powinno~się go nazywać
  \textsc{us-ascii}.

  Powstanie kodu \textsc{ascii} zaczyna~się w~$1961$ roku. Należy sobie
  uświadomić jak wówczas wyglądała świat komputerów. Komputery potrafiły
  zajmować całe pokoje, internet nie był nawet marzeniem, a~przenoszenie
  danych z~jednego komputera polegało na nagraniu danych na kartki papieru
  lub dysk, przeniesienia ich do budynku, gdzie znajduje~się drugi komputer
  i~wgraniu ich do niego.

  I~tu pojawia~się problem. Co jeśli na jednym komputerze zapisywaliśmy
  literę „A” jako liczbę~$1$, ale na drugim komputerze, ktoś przypisał jej
  wartość~$33$? Wbrew pozorom takie sytuacje~się zdarzały dość często
  i~doprowadzały, jak łatwo~się domyślić, użytkowników komputerów do
  rozpaczy.

  Rozwiązanie wydawało~się dość proste. Utwórzmy standard kodowania, którym
  będziemy~się wszyscy posługiwać i~dzięki któremu na każdym komputerze
  literze~„A”, będzie odpowiadała ta sama liczba.

\end{frame}
% ##################





% ##################
\begin{frame}
  \frametitle{Historia kodu ASCII}


  Wprowadzanie tej idei w~życie rozpoczął pracujący wówczas dla firmy
  \textsc{ibm} \colorhref{https://en.wikipedia.org/wiki/Bob\_Bemer}{Bob
    Bemer}, który w~maju $1961$~roku złożył projekt takiego standardu
  do organizacji, która dziś nosi nazwę
  \colorhref{https://en.wikipedia.org/wiki/American\_National\_Standards\_Institute}
  {American National Standards Institute}.

  W~tamtych czasach pamięć była droga, każdy bit~się liczył, zaś bajt liczący
  osiem bitów nie był jeszcze powszechnie przyjętym standardem. Po ustaleniu,
  że~potrzeba im standardu dla około $100$ znaków, wzięto najmniejszą
  liczbę bitów, która zapewnia taką liczbę stanów. Liczbą tą jest~$7$,
  bo~$2^{ \HorSpaceTwo 7 } = 128$. Tym samy kodowanie \textsc{ascii} zawiera
  $128$ znaków, przypisane liczbom od~$0$ do~$127$ (proszę zauważyć,
  że~tych liczb jest dokładnie~$128$).

  Pełną tabelę znaków \textsc{ascii} można znaleźć choćby
  \colorhref{https://www.ascii-code.com/}{tutaj}. Pierwsze $32$ znaki tego
  kodu to tzw. znaki kontrolne (ang. \textit{control characters}). Nie są to
  znaki w~sensie języka codziennego, ile raczej symbole służące do kontroli
  pracy komputera.

\end{frame}
% ##################





% ##################
\begin{frame}
  \frametitle{Struktura kodu ASCII}


  Liczbie $0$ przypisany jest symbol null, służący do~informowania komputera,
  że~coś uległo zakończeniu. Opowiemy sobie o~nim więcej, gdy będziemy
  mówić o~stringach (napisach) w~języku~C. Symbol o~numerze~$2$ informuje
  komputer, że~właśnie rozpoczyna~się pewien tekst, symbol o~numerze~$3$,
  że~ten tekst~się właśnie kończy, symbol~$4$ oznacza zakończenie przesyłu
  danych,~etc. Nie muszą Państwo tego pamiętać, to jest ważne tylko by
  wyjaśnić, czemu \textsc{ascii} ma taki kształt, a~nie inny. Musiał on
  bowiem wziąć pod uwagę wiele zagadnień technicznych, przed którymi
  stanęli jego twórcy w~latach $60$-tych XX~wieku.

  Ostatecznie, literom alfabetu angielskiego przyporządkowano następujące
  liczby. $\text{A} \to 65$, $\text{B} \to 66$, $\text{C} \to 67$, \ldots,
  $\text{Z} \to 90$, $\text{a} \to 97$, $\text{b} \to 98$, $\text{c} \to 99$,
  \ldots, $\text{z} \to 122$. Oprócz tego znakowi przecinka przyporządkowano
  liczbę~$44$, kropce $46$, średnikowi $59$,~etc.

\end{frame}
% ##################






% ##################
\begin{frame}
  \frametitle{Struktura kodu ASCII}


  Ponieważ \textsc{ascii} stworzyli amerykanie, więc zawiera tylko litery
  alfabetu angielskiego. W~jaki sposób kodowane są dziś litery takie jak
  „ą”, czy „ć” to temat na zupełnie inny wykład.

  Nie będziemy od Państwa wymagać znajomości tego wszystkiego. Nasza opowieść
  o~\textsc{ascii} zawiera pewne uproszczenia, gdyż jej główny celem było
  dostarczenie Państwu informacji, pozwalające zrozumieć pewne cechy
  języka~C, które bez tego mogłyby~się wydawać, zupełnie bezsensowne.
  Teraz zaś możemy przejść bezpośrednio do kodowania w~języku~C.

\end{frame}
% ##################










% ######################################
\section{Systemy dziesiętny, binarny i~szesnastkowy}
% ######################################


% ##################
\begin{frame}
  \frametitle{Struktura kodu ASCII}


  \begin{center}

    \begin{tabular}{|c|c|c|c|}
      \hline
      Dec & Bin & Oct & Hex \\
      \hline
      \hphantom{0}$0$ & $0\btext 0$\hphantom{000}
                & $0\otext 0$\hphantom{0} & $0\xtext 0$ \\
      \hphantom{0}$1$ & $0\btext 1$\hphantom{000}
                & $0\otext 1$\hphantom{0} & $0\xtext 1$ \\
      \hphantom{0}$2$ & $0\btext 10$\hphantom{00}
                & $0\otext 2$\hphantom{0} & $0\xtext 2$ \\
      \hphantom{0}$3$ & $0\btext 11$\hphantom{00}
                & $0\otext 3$\hphantom{0} & $0\xtext 3$ \\
      \hphantom{0}$4$ & $0\btext 100$\hphantom{0}
                & $0\otext 4$\hphantom{0} & $0\xtext 4$ \\
      \hphantom{0}$5$ & $0\btext 101$\hphantom{0}
                & $0\otext 5$\hphantom{0} & $0\xtext 5$ \\
      \hphantom{0}$6$ & $0\btext 110$\hphantom{0}
                & $0\otext 6$\hphantom{0} & $0\xtext 6$ \\
      \hphantom{0}$7$ & $0\btext 111$\hphantom{0}
                & $0\otext 7$\hphantom{0} & $0\xtext 7$ \\
      \hphantom{0}$8$ & $0\btext 1000$ & $0\otext 10$ & $0\xtext 8$ \\
      \hphantom{0}$9$ & $0\btext 1001$ & $0\otext 11$ & $0\xtext 9$ \\
      $10$ & $0\btext 1010$ & $0\otext 12$ & $0\xtext \Atext$ \\
      $11$ & $0\btext 1011$ & $0\otext 13$ & $0\xtext \Btext$ \\
      $12$ & $0\btext 1100$ & $0\otext 14$ & $0\xtext \Ctext$ \\
      \hline
    \end{tabular}

  \end{center}

\end{frame}
% ##################





% ##################
\begin{frame}
  \frametitle{Struktura kodu ASCII}


  \begin{center}

    \begin{tabular}{|c|c|c|c|}
      \hline
      Dec & Bin & Oct & Hex \\
      \hline
      $13$ & $0\btext 1101$\hphantom{0} & $0\otext 15$ & $0\xtext \Dtext$ \\
      $14$ & $0\btext 1110$\hphantom{0} & $0\otext 16$ & $0\xtext \Etext$ \\
      $15$ & $0\btext 1111$\hphantom{0} & $0\otext 17$ & $0\xtext \Ftext$ \\
      $16$ & $0\btext 1 0000$ & $0\otext 20$ & $0\xtext 10$ \\
      $17$ & $0\btext 1 0001$ & $0\otext 21$ & $0\xtext 11$ \\
      $18$ & $0\btext 1 0010$ & $0\otext 22$ & $0\xtext 12$ \\
      $19$ & $0\btext 1 0011$ & $0\otext 23$ & $0\xtext 13$ \\
      $20$ & $0\btext 1 0100$ & $0\otext 24$ & $0\xtext 14$ \\
      $21$ & $0\btext 1 0101$ & $0\otext 25$ & $0\xtext 15$ \\
      $22$ & $0\btext 1 0110$ & $0\otext 26$ & $0\xtext 16$ \\
      $23$ & $0\btext 1 0111$ & $0\otext 27$ & $0\xtext 17$ \\
      $24$ & $0\btext 1 1000$ & $0\otext 30$ & $0\xtext 18$ \\
      $25$ & $0\btext 1 1001$ & $0\otext 31$ & $0\xtext 19$ \\
      \hline
    \end{tabular}

  \end{center}

\end{frame}
% ##################





% ##################
\begin{frame}
  \frametitle{Struktura kodu ASCII}


  \begin{center}

    \begin{tabular}{|c|c|c|c|}
      \hline
      Dec & Bin & Oct & Hex \\
      \hline
      $26$ & $0\btext 1 1010$\hphantom{0} & $0\otext 32$
                      & $0\xtext 1\Atext$ \\
      $27$ & $0\btext 1 1011$\hphantom{0} & $0\otext 33$
                      & $0\xtext 1\Btext$ \\
      $28$ & $0\btext 1 1100$\hphantom{0} & $0\otext 34$
                      & $0\xtext 1\Ctext$ \\
      $29$ & $0\btext 1 1101$\hphantom{0} & $0\otext 35$
                      & $0\xtext 1\Dtext$ \\
      $30$ & $0\btext 1 1110$\hphantom{0} & $0\otext 36$
                      & $0\xtext 1\Etext$ \\
      $31$ & $0\btext 1 1111$\hphantom{0} & $0\otext 37$
                      & $0\xtext 1\Ftext$ \\
      $32$ & $0\btext 10 0000$ & $0\otext 40$ & $0\xtext 20$ \\
      $33$ & $0\btext 10 0001$ & $0\otext 41$ & $0\xtext 21$ \\
      $34$ & $0\btext 10 0010$ & $0\otext 42$ & $0\xtext 22$ \\
      $35$ & $0\btext 10 0011$ & $0\otext 43$ & $0\xtext 23$ \\
      $36$ & $0\btext 10 0100$ & $0\otext 44$ & $0\xtext 24$ \\
      $37$ & $0\btext 10 0101$ & $0\otext 45$ & $0\xtext 25$ \\
      $38$ & $0\btext 11 0110$ & $0\otext 46$ & $0\xtext 26$ \\
      \hline
    \end{tabular}

  \end{center}

\end{frame}
% ##################





% ##################
\begin{frame}
  \frametitle{Struktura kodu ASCII}


  \begin{center}

    \begin{tabular}{|c|c|c|c|}
      \hline
      Dec & Bin & Oct & Hex \\
      \hline
      $26$ & $0\btext 1 1010$\hphantom{0} & $0\otext 32$
                      & $0\xtext 1\Atext$ \\
      $27$ & $0\btext 1 1011$\hphantom{0} & $0\otext 33$
                      & $0\xtext 1\Btext$ \\
      $28$ & $0\btext 1 1100$\hphantom{0} & $0\otext 34$
                      & $0\xtext 1\Ctext$ \\
      $29$ & $0\btext 1 1101$\hphantom{0} & $0\otext 35$
                      & $0\xtext 1\Dtext$ \\
      $30$ & $0\btext 1 1110$\hphantom{0} & $0\otext 36$
                      & $0\xtext 1\Etext$ \\
      $31$ & $0\btext 1 1111$\hphantom{0} & $0\otext 37$
                      & $0\xtext 1\Ftext$ \\
      $32$ & $0\btext 10 0000$ & $0\otext 40$ & $0\xtext 20$ \\
      $33$ & $0\btext 10 0001$ & $0\otext 41$ & $0\xtext 21$ \\
      $34$ & $0\btext 10 0010$ & $0\otext 42$ & $0\xtext 22$ \\
      $35$ & $0\btext 10 0011$ & $0\otext 43$ & $0\xtext 23$ \\
      $36$ & $0\btext 10 0100$ & $0\otext 44$ & $0\xtext 24$ \\
      $37$ & $0\btext 10 0101$ & $0\otext 45$ & $0\xtext 25$ \\
      $38$ & $0\btext 11 0110$ & $0\otext 46$ & $0\xtext 26$ \\
      \hline
    \end{tabular}

  \end{center}

\end{frame}
% ##################


  miejscach musimy podawać liczby w~zakresie od~$0$ do~$255$. Po porostu
  komputer na zapisanie tej liczby poświęca dokładnie jeden bajt ($8$~bitów).

  System który omówiliśmy nie pozwala zapisywać choć liczb całkowitych
  ujemnych, musi więc zostać odpowiednio rozszerzony. Jednak my nie będziemy
  wnikać w~tego typu zagadnienia, gdyż to co powiedzieliśmy, jest na nasze
  potrzeby wystarczające.


    Jak dobrze wszyscy wiemy, jednostkę która może przechowywać wartość
  $0$ albo $1$ nazywamy \textbf{bitem}. Bity grupujemy dziś
  w~\textbf{bajty},
  jeden bajt to $2^{ \HorSpaceTwo 3 } == 8$ bitów. W~tym, że~bajt zawiera
  $2^{ \HorSpaceTwo 3 }$ bitów widać upodobanie informatyki do potęgi
  liczby~$2$. Skąd jednak jest to upodobanie.

  Wydaje~się, że~podstawą tego są kwestie techniczne. Zbudowanie urządzenia,
  które może przechowywać dwie różne wartości ($0$ lub $1$), jest prostsze,
  niż takiego, które przechowywałoby trzy wartości ($0$, $1$ lub $2$).
  Warto powiedzieć o~tym kilka słów więcej.

  Zadajmy sobie jeszcze pytanie, ile liczb jesteśmy w~stanie zapisać


  Te przykłady powinny zilustrować następujący fakt. Mając dostępne
  $n$~bitów możemy za ich pomocą zapisać $2^{ \HorSpaceTwo n }$ stanów. Jest
  to fakt ogólny, którego prosto dowodzi~się matematycznie, jednak
  zajmowanie~się matematycznymi podstawami informatyki, nie należy do zakresu
  tych zajęć. Naszym celem jest pomóc wyrobić sobie Państwu odpowiednią
  intuicję, pozostawiając innym przedmiotom dokładniejsze omówienie tych
  spraw.

  Zwróćmy uwagę na jeszcze jeden fakt. W~tym systemie notacji liczb
  naturalnych, mając do dyspozycji $n$~bitów jeden stan zużywamy na zapisanie
  liczby $0 == 0\text{b}00\ldots0$, pozostałe $2^{ \HorSpaceTwo n - 1 }$ zużywamy
  na zapisanie kolejny liczb naturalnych. Stąd, w~tym konkretnym
  systemie, największą liczbą jaką możemy zapisać na $n$~bitach
  jest $2^{ \HorSpaceTwo n } - 1$.

    Pozostając na bardzo ogólnym poziomie, w~XX wieku eksperymentowano
  z~układami elektronicznymi, których podstawowa jednostka pamięci, była
  w~stanie przechowywać jedną z~trzech wartości ($0$, $1$ lub $2$), jedną
  z~czterech ($0$, $1$, $2$ lub $3$),~etc. Jednostki te odpowiadały temu,
  czemu w~dzisiejszym komputerze odpowiada bit. Gdyby ludzkość
  zdecydowała~się budować komputery w systemie trójkowym, pewnie dzisiaj
  operowalibyśmy „bajtami” o~długości $3^{ 2 } == 9$ „bitów”.

  Różne względy, które są bardziej odpowiednie na bardzo obszerny wykład
  z~elektroniki, sprawiły, iż~zdecydowaliśmy~się oprzeć pamięć komputera
  o~jednostki dwustanowe, czyli właśnie dobrze nam znane bity. Czy tak już
  pozostanie? Przewidywanie przyszłości jest trudne.








% ######################################
\section{Wprowadzenie do liczb pseudolosowych}
% ######################################


% ##################
\begin{frame}
  \frametitle{Czemu potrzebujemy liczb pseudolosowych?}


  Zacznijmy od dyskusji problemu pozornie odległego od informatyki. Mamy
  fabrykę, która produkuje maszynki do golenia. Wyprodukowaliśmy milion
  maszynek, wzięliśmy jedną do ręki, ona zaś złamała~się od razu. Jeśli
  wyślemy takie maszynki do sklepów, bo klienci przestaną je kupować, a~my
  tego nie chcemy. Stajemy jednak przed pytaniem, ile maszynek z~miliona
  jest tak niskiej jakości? Jeśli jest ich $900 \, 000$ to lepiej wszystkie
  policzyć na straty, bo sprzedaż ich tylko zniszczy nam reputację. Co jeśli
  jednak tylko $1000$ jest w~tak złym stanie i~my przypadkowo wybraliśmy
  jedną z~nich? Wtedy oczywiście należy wysłać cały milion do sklepów, gdyż
  $1000$ klientom którzy mieli pecha, po prostu zwrócimy maszynki na koszt
  firmy.

  Jak jednak stwierdzić, ile maszynek jest dobrych, a~ile złych? Nie
  sprawdzimy każdej z~miliona, bo to za dużo kosztuje. Możemy jednak
  sprawdzić spokojnie $100$. Tylko które $100$?

\end{frame}
% ##################





% ##################
\begin{frame}
  \frametitle{Czemu potrzebujemy liczb pseudolosowych?}


  Jeśli sprawdzimy ostatnie $100$ wyprodukowanych, to może się okazać,
  że~maszyna zepsuła~się pod koniec, ostatni $1000$ maszynek jest przez to
  do utylizacji, ale wszystkie inne są dobre. Jeśli więc będziemy mieli
  pecha, to wyciągniemy błędny wniosek z~jakości $100$ maszynek i~zmarnujemy
  dobre $999 \, 000$.

  To może wziąć $100$ pierwszych maszynek? Ale co jeśli maszyna zepsuła~się
  po wyprodukowaniu pierwszego tysiąca, wtedy $100$ pierwszych maszynek
  jest dobra, a~my wyślemy do sklepów $999 \, 000$ zepsutych, co zrujnuje
  nam reputację. To może wziąć $100$ maszynek ze środka?

  Dziedziny matematyki jakimi są rachunek prawdopodobieństwa i~statystyka,
  podpowiadają nam jakie są optymalne strategie radzenia sobie z~takimi
  problemami. Ponieważ wątpimy, że~Państwo chcą słuchać o~tych sprawach,
  wspominamy o~tym tylko dlatego, że~stanowi to motywację, dla tego
  co będziemy robić.

\end{frame}
% ##################





% ##################
\begin{frame}
  \frametitle{Czemu potrzebujemy liczb pseudolosowych?}


  Jedną z~dobrych strategii poradzenia sobie z~tym problemem, jest
  ponumerowanie wszystkich maszynek od~$1$ do~$1 \, 000 \, 000$, wylosowanie
  $100$ liczb z~tego zakresu w~sposób losowy, następnie przetestowanie
  tak wylosowanych maszynek. W~zależności od tego jak wiele z~tych maszynek
  było dobrej jakości, a~ile zepsutych, podejmiemy decyzję co zrobić
  z~całą liczącą $1 \, 000 \, 000$ sztuk produkcją. Oczywiście, może~się
  zdarzyć, iż~mamy wielkiego pecha i~wyciągniemy zupełnie błędny wniosek,
  ale użyta metoda pozwala bardzo mocno ograniczyć wystąpienie takich
  sytuacji.

  Dzisiaj takie analizy testowania jakości potrafią być już bardzo złożone
  i~powyższy przykład nie powinien być traktowany jako nic więcej, niż
  prosta ilustracja problemu. Nas zaś zostawia z~następującym pytaniem: skąd
  wziąć $100$ liczb losowy z~zadanego przedziału?

\end{frame}
% ##################





% ##################
\begin{frame}
  \frametitle{Skąd wziąć liczby (pseudo)losowe?}


  W~dzisiejszych czasach zwykle wtedy siadamy do komputera i~prosimy go
  o~wygenerowanie nam $100$~odpowiednich liczb. I~tutaj wracamy do
  informatyki.

  Okazuje~się, że~znalezienie liczb naprawdę losowy jest niezwykle trudne.
  Teoria prawdopodobieństwa i~statystyka pozwalają nam znaleźć prawie
  idealne rozwiązania wielu praktycznych problemów, jeśli tylko umiem
  znaleźć prawdziwie losowy zestaw liczb. W~poprzednim przykładzie były to
  numery wyprodukowanych maszynek. Włożono więc bardzo wiele wysiłku
  w~szukanie liczb (pseudo)losowych oraz testowanie, czy dany ciąg liczb
  jest odpowiednio dobry.

  Dzisiejsze wymagania nauki, techniki i~gospodarki sprawiły, że~nasze metody
  badania czy coś jest losowe, stały~się niezwykle subtelne i~wiele źródeł
  liczb losowych musiało zostać odrzuconych jako nie dość dobre.

\end{frame}
% ##################





% ##################
\begin{frame}
  \frametitle{Skąd wziąć liczby (pseudo)losowe?}


  Przykładowo biorę monetę i~wyrzucam ciąg $\text{O}$, $\text{O}$,
  $\text{O}$, $\text{R}$ ($\text{O}$~-- orzeł, $\text{R}$~-- reszka)
  i~przypisuję orłowi wartość~$1$, a~reszce~$0$, dostaje więc liczbę
  $0\text{b}1110 == 14$. Jednak moneta nie jest nigdy idealnie wyważona,
  więc rzucając nią, będziemy dostawali, przykładowo, trochę więcej orłów,
  niż reszek. Nasze dzisiejsze wymagania odnośnie jakości liczb losowych
  są tak wysokie, że~tego typu odstępstwa od prawdziwej losowości, już
  skreślają daną metodę dla nas.

  Również tego przykładu nie należy traktować specjalnie poważnie. Od jego
  zgodności z~rzeczywistością ważniejsza jest ogólna idea, która za nim
  stoi.

  Mam nadzieję, że~to Państwa przynajmniej trochę przekonało, że~liczby
  losowe są przydatne w~wielu praktycznych problemach, również
  w~informatyce. Jednak przykład z~fabryką maszynek jest prostszy do
  zrozumienia, niż problemy napotykane w~informatyce.

\end{frame}
% ##################





% ##################
\begin{frame}
  \frametitle{Skąd wziąć liczby (pseudo)losowe?}


  Jeśli nie umiemy znaleźć liczb naprawdę losowych, to może uda nam~się
  znaleźć liczby, które są wystarczająco podobne do losowych? Takie
  właśnie liczby, które nie są naprawdę losowe, ale są do nich
  wystarczająco podobne nazywamy liczbami
  \colorhref{https://en.wikipedia.org/wiki/Pseudorandomness}
  {pseudolosowymi}. Generowanie takich liczb to wciąż otwarta dziedzina
  badań, ale najbardziej podstawowe idea nie są trudne do zrozumienia.

  Zacznijmy od znalezienia funkcji, która zmienia~się w~bardzo
  „nieprzewidywalny” sposób. Przykład takiej funkcji napisanej w~języku~C,
  można znaleźć w~pliku ????

  Zaprezentowana tam funkcja \texttt{generatePseudoRandomNumber()} jest
  funkcją deterministyczną, czyli \texttt{generatePseudoRandomNumber(0)}
  zawsze zwraca liczbę $33$, \texttt{generatePseudoRandomNumber(1)}
  zwraca $308$, \texttt{generatePseudoRandomNumber(2)} zwraca $34$, etc.
  Jednak powiedzenie nawet w~przybliżeniu, ile będzie wynosiła wartość
  tej funkcji na wybrany argumencie, nie jest zadaniem prostym.

\end{frame}
% ##################





% ##################
\begin{frame}
  \frametitle{Skąd wziąć liczby (pseudo)losowe?}


  Mają już taką „nieprzewidywalną” funkcję, wybieramy liczbę zwaną
  „ziarnem” (ang.~\textit{seed}). Wybór ziarna jest zupełnie dowolny,
  możemy je na przykład otrzymać rzucają kostką. Następnie obliczamy
  wartość \texttt{generatePseudoRandomNumber(ziarno)}. Myśl za tym stoi
  jest mniej więcej taka. Jeśli mam „trochę losową” liczbę jaką jest ziarno
  i~wrzucę ją do „nieprzewidywalnej” funkcji, to rezultat zwrócony przez tę
  funkcję będzie „bardziej losowy”, niż to od czego zacząłem.

  Istnieje wielka literatura na temat tego, co to znaczy być mniej lub
  bardziej losowym i~jestem przekonany, że~nikt z~Państwa nie chce o~tym
  słuchać. Na tym kursie zupełnie wystarczy intuicyjne zrozumienie tego
  zagadnienia.

\end{frame}
% ##################





% ##################
\begin{frame}
  \frametitle{Skąd wziąć liczby (pseudo)losowe?}


  Wróćmy jeszcze na chwilę do naszej funkcji generującej pseudolosowe liczby
  \texttt{generatePseudoRandomNumber(ziarno)}. Funkcje tego typu
  noszą nazwę \textbf{generatorów liczb pseudolosowych}, choć niewątpliwie
  ta konkretna funkcja jest wyjątkowo kiepskim generatorem takich liczb.
  Jestem ze sposób poprawienia działania takich funkcji jest wywołanie
  jej na samej sobie, w~sposób pokazany w~pliku ???.

  To jakiego generator liczb pseudolosowych należy użyć, zależy w~dużej
  mierze od~tego, jak dobrych liczb tego typu potrzebujemy by~rozwiązać
  nasz problem? Jako, że~to jest tylko kurs \textit{Podstaw informatyki},
  chcieliśmy jedynie zarysować przed Państwem ten problem oraz pokazać,
  jakie narzędzia daje nam w~tej kwestii język~C.

\end{frame}
% ##################




% % ############################
% \jagiellonianendslide{Czy są jakieś pytania do tej części?}
% % ############################































% ####################################################################
% ####################################################################
% Bibliography

\printbibliography





% ############################
% End of the document

\end{document}
